% !TeX program = xelatex
\documentclass[12pt,a4paper]{article}

% =========================================================
% ЯЗЫК И ШРИФТЫ
% =========================================================

\usepackage{fontspec}
\usepackage{polyglossia}
\usepackage{dsfont}

\setdefaultlanguage{russian}
\setotherlanguage{english}

\setmainfont{CMU Serif}[ItalicFont=cmunti.otf]
\setsansfont{CMU Sans Serif}
\setmonofont{CMU Typewriter Text}


% =========================================================
% МАТЕМАТИКА
% =========================================================
\usepackage{amsmath}
\usepackage{amssymb}
\usepackage{amsfonts}
\usepackage{amsthm}
\usepackage{mathtools}
\usepackage{bm}
\usepackage{mathrsfs}

% =========================================================
% СТРАНИЦА
% =========================================================
\usepackage[a4paper, margin=2cm]{geometry}

% =========================================================
% ЦВЕТА И РАМКИ
% =========================================================
\usepackage{xcolor}
\usepackage[most]{tcolorbox}

% =========================================================
% ГРАФИКА
% =========================================================
\usepackage{graphicx}

% =========================================================
% СПИСКИ
% =========================================================
\usepackage{enumitem}

% =========================================================
% ЗАГОЛОВКИ
% =========================================================
\usepackage{titlesec}

% =========================================================
% КОЛОНТИТУЛЫ
% =========================================================
\usepackage{fancyhdr}

% =========================================================
% ТАБЛИЦЫ
% =========================================================
\usepackage{booktabs}
\usepackage{array}

% =========================================================
% ССЫЛКИ
% =========================================================
\usepackage{hyperref}

\hypersetup{
    colorlinks=true,
    linkcolor=blue,
    urlcolor=blue,
    citecolor=blue
}

% =========================================================
% НАСТРОЙКИ
% =========================================================
\setlength{\parindent}{0pt}
\setlength{\parskip}{6pt}

% -------------------------
% Определение
% -------------------------
\newtcolorbox{definition}[1][]{
    breakable,
    colback=red!5,
    colframe=red!75!black,
    fonttitle=\bfseries,
    title=Определение,
    #1
}

% -------------------------
% Теорема
% -------------------------
\newtcolorbox{theorem}[2][]{
    breakable,
    colback=green!5,
    colframe=green!50!black,
    fonttitle=\bfseries,
    title=Теорема #2,
    #1
}

% -------------------------
% Лемма
% -------------------------
\newtcolorbox{lemma}[2][]{
    breakable,
    colback=orange!5,
    colframe=orange!70!black,
    fonttitle=\bfseries,
    title=Лемма #2,
    #1
}

% -------------------------
% Следствие
% -------------------------
\newtcolorbox{corollary}[1][]{
    breakable,
    colback=green!5,
    colframe=green!35!black,
    fonttitle=\bfseries,
    title=Следствие,
    #1
}

% -------------------------
% Доказательство
% -------------------------
\newtcolorbox{proofbox}[2][]{
    breakable,
    colback=blue!5,
    colframe=blue!50,
    fonttitle=\bfseries,
    title=Доказательство #2,
    #1
}

% -------------------------
% Замечание
% -------------------------
\newtcolorbox{note}[1][]{
    breakable,
    colback=purple!5,
    colframe=purple!80!black,
    fonttitle=\bfseries,
    title=Замечание,
    #1
}

% -------------------------
% Свойства
% -------------------------
\newtcolorbox{props}[2][]{
    breakable,
    colback=green!5,
    colframe=green!70!black,
    fonttitle=\bfseries,
    title=Свойства #2,
    #1
}

% -------------------------
% Примеры
% -------------------------
\newtcolorbox{example}[1][]{
    breakable,
    colback=brown!5,
    colframe=brown!40!black,
    fonttitle=\bfseries,
    title=Примеры,
    #1
}

% =========================================================
% ДОКУМЕНТ
% =========================================================
\begin{document}

\tableofcontents

\newpage

\section{Вероятностное пространство}

\subsection{\texorpdfstring{$\sigma$-алгебра}{sigma-алгебра}}

\begin{definition}
\textbf{$\sigma$-алгебра} $\mathcal{A}$ на множестве $\Omega$ — непустое семейство подмножеств $\Omega$, такое что:

\[
A\in\mathcal{A}
\Longrightarrow
\overline{A}\in\mathcal{A},
\]

и для любой последовательности множеств $(A_k)_{k=1}^{\infty}$ из $\mathcal{A}$

\[
\bigcup_{k=1}^{\infty}A_k\in\mathcal{A}
\]

\end{definition}

\begin{props}{$\sigma$-алгебры}

1) \textit{Замкнутость относительно пересечения.}

Для любой последовательности $(A_k)_{k=1}^{\infty}$ из $\mathcal{A}$

\[
\bigcap_{k=1}^{\infty}A_k
=
\overline{
\bigcup_{k=1}^{\infty}\overline{A_k}
}
\in\mathcal{A}
\]

\medskip

2) Если $A,B\in\mathcal{A}$, то

\[
A\setminus B
=
A\cap\overline{B}
\in\mathcal{A}
\]

\medskip

3) Для любого $A\in\mathcal{A}$

\[
\Omega=A\cup\overline{A}\in\mathcal{A}
\]

\medskip

4)

\[
\emptyset=\overline{\Omega}\in\mathcal{A}
\]

\end{props}


\begin{example}

1) Полная $\sigma$-алгебра:

\[
2^\Omega
\]

\medskip

2) Тривиальная $\sigma$-алгебра:

\[
\{\emptyset,\Omega\}
\]

\end{example}


\subsection{Вероятность}

\begin{definition}

Пусть $\Omega$ — множество элементарных исходов, а $\mathcal{A}$ — $\sigma$-алгебра.

Отображение

\[
P:\mathcal{A}\to\mathbb{R}
\]

называется \textbf{вероятностью}, если:

\[
P(A)\geq0;
\]

\[
P(\Omega)=1;
\]

и для попарно несовместных событий $A_k\in\mathcal{A}$

\[
P\left(
\bigsqcup_{k=1}^{\infty}A_k
\right)
=
\sum_{k=1}^{\infty}P(A_k)
\]

Тройка

\[
(\Omega,\mathcal{A},P)
\]

называется \textbf{вероятностным пространством}.

\end{definition}


\begin{example}

1) Общая дискретная модель:

\[
\Omega=\{\omega_k\}_{k=1}^{N},
\qquad N\leq\infty,
\]

\[
\mathcal{A}=2^\Omega,
\]

\[
P(\{\omega_k\})
=
P(\omega_k)
=
p_k\in[0,1],
\]

\[
\sum_{k=1}^{N}p_k=1
\]

Если

\[
A=\{\omega_{i_j}\}_{j=1}^{m},
\]

то

\[
P(A)
=
\sum_{j=1}^{m}P(\omega_{i_j})
=
\sum_{j=1}^{m}p_{i_j}
\]


\medskip

2) Классическая вероятностная модель:

\[
\Omega=\{\omega_1,\ldots,\omega_N\},
\]

\[
P(\omega_k)=\frac1N
\]

Тогда

\[
P(A)=\frac{|A|}{N}
\]


\medskip

3) Последовательность независимых испытаний с двумя исходами:

\[
\Omega=\{0,1\}^{\mathbb{N}}
\]

Если вероятность появления $1$ равна $p$, а вероятность появления $0$ равна

\[
q=1-p,
\]

то для последовательности, содержащей $k$ единиц и $n-k$ нулей:

\[
P(0^{\,n-k}1^k)
=
q^{\,n-k}p^k
\]


\medskip

4) Геометрическая вероятность:

\[
\Omega\subset\mathbb{R}^n,
\qquad
0<\operatorname{Vol}\Omega<\infty
\]

Система событий:

\[
\mathcal{A}
=
\left\{
A\subset\Omega
\mid
\operatorname{Vol}A\in
[0,\operatorname{Vol}\Omega]
\right\}
\]

Тогда

\[
P(A)
=
\frac{\operatorname{Vol}A}
{\operatorname{Vol}\Omega}
\]


\begin{note}
$\mathcal{A}$ не всегда равно $2^\Omega$, так как существуют множества, которые нельзя измерить.
\end{note}

\end{example}

\subsection{Алгебра событий}

\begin{definition}

\textbf{Алгебра событий} $\mathbf{A}$ — непустое семейство подмножеств $\Omega$, такое что:

\[
A\in\mathbf{A}
\Longrightarrow
\overline{A}\in\mathbf{A},
\]

и

\[
A,B\in\mathbf{A}
\Longrightarrow
A\cup B\in\mathbf{A}
\]

\end{definition}


\begin{note}

1)

\[
\mathbf{A}\subset2^\Omega
\]

То есть существует $\sigma$-алгебра, покрывающая $\mathbf{A}$.


\medskip

2)

\[
\sigma(\mathbf{A})
=
\bigcap\mathcal{A}
\]

— минимальная $\sigma$-алгебра.

\end{note}


\begin{definition}

\textbf{Борелевская $\sigma$-алгебра} $\mathcal{B}$ — $\sigma$-алгебра, содержащая интервалы всех видов.

В частности, при

\[
-\infty<a<b<+\infty
\]

она содержит множества вида

\[
[a,b]
=
\overline{
(-\infty,a)\cup(b,+\infty)
}
\]

Также

\[
[a,b)
=
\overline{
(-\infty,a)\cup[b,c]\cup(c,+\infty)
}
\]

Кроме того,

\[
\{a\}
=
\bigcap_{n=1}^{\infty}
\left[a,a+\frac1n\right]
\]

\end{definition}


\begin{theorem}{о продолжении меры}

Пусть $\Omega$ — множество, $\mathcal{A}$ — алгебра событий и

\[
P:\mathcal{A}\to\mathbb{R}
\]

удовлетворяет условиям:

\[
P(A)\geq0,
\]

\[
P(\Omega)=1,
\]

и для любой последовательности попарно несовместных множеств $A_k\in\mathcal{A}$, такой что

\[
\bigsqcup_{k=1}^{\infty}A_k\in\mathcal{A},
\]

выполнено:

\[
P\left(
\bigsqcup_{k=1}^{\infty}A_k
\right)
=
\sum_{k=1}^{\infty}P(A_k)
\]

Тогда существует единственная вероятность

\[
Q:\sigma(\mathcal{A})\to\mathbb{R},
\]

продолжающая $P$:

\[
Q|_{\mathcal{A}}=P
\]

\end{theorem}


\begin{props}{вероятности}

1) Если $A\subset B$, то

\[
P(A)\leq P(B)
\]


2) Для любого события $A$

\[
P(A)\leq1
\]


3) Для любого события $A$

\[
P(A)=1-P(\overline{A})
\]


4) Для любых событий $A$ и $B$

\[
P(A\cup B)
=
P(A)+P(B)-P(A\cap B)
\]


5) $P(\emptyset)=0$



6) Для любой последовательности событий $(A_k)_{k=1}^{\infty}$

\[
P\left(
\bigcup_{k=1}^{\infty}A_k
\right)
\leq
\sum_{k=1}^{\infty}P(A_k)
\]

\end{props}

\begin{proofbox}{свойства 1}

Имеем дизъюнктное разбиение

\[
B=A\sqcup(B\setminus A)
\]

Следовательно,

\[
P(B)
=
P(A)+P(B\setminus A)
\]

Так как

\[
P(B\setminus A)\geq0,
\]

то

\[
P(B)\geq P(A)
\]

\end{proofbox}

\begin{proofbox}{свойства 2}

Так как

\[
A\subset\Omega,
\]

по первому свойству получаем

\[
P(A)\leq P(\Omega)
\]

Но

\[
P(\Omega)=1,
\]

следовательно,

\[
P(A)\leq1
\]

\end{proofbox}

\begin{proofbox}{свойства 3}

Так как

\[
\Omega=A\sqcup\overline{A},
\]

то по свойству аддитивности вероятности

\[
P(\Omega)
=
P(A)+P(\overline{A})
\]

Так как

\[
P(\Omega)=1,
\]

получаем

\[
1=P(A)+P(\overline{A}),
\]

откуда

\[
P(A)=1-P(\overline{A})
\]

\end{proofbox}

\begin{proofbox}{свойства 4}

Разложим $B$ в дизъюнктное объединение:

\[
B=(B\setminus(A\cap B))\sqcup(A\cap B)
\]

Тогда

\[
P(B)
=
P(B\setminus(A\cap B))
+
P(A\cap B)
\]

Также

\[
A\cup B
=
A\sqcup(B\setminus(A\cap B)),
\]

поэтому

\[
P(A\cup B)
=
P(A)+P(B\setminus(A\cap B))
\]

Следовательно,

\[
P(B\setminus(A\cap B))
=
P(B)-P(A\cap B)
\]

Подставляя, получаем

\[
P(A\cup B)
=
P(A)+P(B)-P(A\cap B)
\]

\end{proofbox}

\begin{proofbox}{свойства 5}

Так как

\[
\emptyset=\Omega\setminus\Omega,
\]

то

\[
\Omega=\Omega\sqcup\emptyset
\]

Следовательно,

\[
P(\Omega)
=
P(\Omega)+P(\emptyset)
\]

Так как

\[
P(\Omega)=1,
\]

получаем

\[
P(\emptyset)=0
\]

\end{proofbox}

\begin{proofbox}{свойства 6}

Положим

\[
B_1=A_1,
\]

а для $k\geq2$

\[
B_k
=
A_k\setminus
\bigcup_{i=1}^{k-1}A_i
\]

Тогда множества $B_k$ попарно несовместны и

\[
\bigcup_{k=1}^{\infty}A_k
=
\bigsqcup_{k=1}^{\infty}B_k
\]

Кроме того,

\[
B_k\subset A_k
\]

Поэтому

\[
P\left(
\bigcup_{k=1}^{\infty}A_k
\right)
=
\sum_{k=1}^{\infty}P(B_k)
\]

Так как

\[
P(B_k)\leq P(A_k),
\]

то

\[
P\left(
\bigcup_{k=1}^{\infty}A_k
\right)
\leq
\sum_{k=1}^{\infty}P(A_k)
\]

\end{proofbox}

\begin{theorem}{- Формула включений и исключений}
Пусть $A_1,\dots,A_n$ --- события. Тогда
\[
\begin{aligned}
P(A_1\cup\dots\cup A_n)
={}&\sum_{k=1}^{n}P(A_k)
-\sum_{k<l}P(A_k\cap A_l)\\
&+\sum_{k<l<m}P(A_k\cap A_l\cap A_m)
+\dots\\
&+(-1)^{n-1}P(A_1\cap\dots\cap A_n)
\end{aligned}
\]
\end{theorem}

\begin{theorem}{- Свойства условной вероятности}
Следующие условия равносильны:

\begin{enumerate}
    \item $P$ конечно-аддитивно.
    \item $P$ конечно аддитивно и для
    \[
    B_1\supseteq B_2\supseteq\dots\supseteq B_n,\qquad
    B=\bigcap B_k
    \]
    выполняется
    \[
    P(B_k)\to P(B)
    \]
    (непрерывность сверху).

    \item $P$ конечно аддитивно и для
    \[
    A_1\subseteq A_2\subseteq\dots,\qquad
    A=\bigcup A_k
    \]
    выполняется
    \[
    P(A_k)\to P(A)
    \]
    (непрерывность снизу).
\end{enumerate}
\end{theorem}

\begin{proofbox}{непрерывности снизу из непрерывности сверху}
\textbf{Непрерывность сверху $\Rightarrow$ непрерывность снизу.}

Пусть
\[
A_k\subseteq A_{k+1},\qquad A=\bigcup A_k
\]
Рассмотрим дополнения:
\[
B_k=\overline{A_k},\qquad B=\overline{A}
\]
Тогда
\[
B_k\supseteq B_{k+1},\qquad B=\bigcap B_k
\]

По непрерывности сверху:
\[
P(B_k)\to P(B)
\]

Так как
\[
P(B_k)=1-P(A_k),\qquad P(B)=1-P(A),
\]
получаем
\[
P(A_k)\to P(A)
\]
\end{proofbox}

\begin{proofbox}{непрерывности сверху из непрерывности снизу}
\textbf{Непрерывность снизу $\Rightarrow$ непрерывность сверху.}

Пусть
\[
B_k\supseteq B_{k+1},\qquad B=\bigcap B_k
\]

Рассмотрим множества
\[
C_k=B_1\setminus B_k
\]
Тогда
\[
C_1\subseteq C_2\subseteq\dots,
\qquad
C=\bigcup C_k=B_1\setminus B
\]

По непрерывности снизу:
\[
P(C_k)\to P(C)
\]

Так как
\[
P(C_k)=P(B_1)-P(B_k),
\]
то
\[
P(B_k)\to P(B)
\]
\end{proofbox}


\begin{proofbox}{счётной аддитивности из непрерывности сверху}
\textbf{Непрерывность сверху $\Rightarrow$ счётная аддитивность.}

Пусть
\[
B_k:B_i\cap B_j=\varnothing,\quad i\neq j
\]

Обозначим
\[
C_1=B_1\setminus B_2,\qquad
C_2=B_2\setminus B_3,\qquad
\dots
\]
Тогда
\[
B_j=\bigcup_{k\leq j}C_k\cup B
\]

Переходя к пределу:
\[
P(B_j)\to P(B)+\sum_{k=1}^{\infty}P(C_k)
\]

Так как остаток стремится к нулю, получаем счётную аддитивность.
\end{proofbox}


\begin{proofbox}{непрерывности сверху из счётной аддитивности}
\textbf{Счётная аддитивность $\Rightarrow$ непрерывность сверху.}

Пусть
\[
D_k\downarrow\varnothing
\]
Тогда
\[
\bigcup_{k=1}^{\infty}D_k
=
D_1\cup\bigcup_{k=1}^{\infty}(D_k\setminus D_{k+1})
\]

По счётной аддитивности:
\[
P(D_1)=P\left(\bigcap_{k=1}^{\infty}D_k\right)
+
\sum_{k=1}^{\infty}P(D_k\setminus D_{k+1})
\]

Следовательно,
\[
P(D_k)\to P\left(\bigcap_{k=1}^{\infty}D_k\right)
\]
\end{proofbox}


\subsection{Пример неизмеримого множества}

\begin{example}
Рассмотрим множество
\[
\Omega=\{-1,1\}
\]

Определим отношение эквивалентности:
\[
a\sim b\Longleftrightarrow a-b\in\mathbb{Q}
\]

Выберем по одному представителю из каждого класса эквивалентности.

Обозначим полученное множество через $E_0$.

Не существует такой меры, которая каждому ограниченному множеству присваивает неограниченную меру.

\[
A\subset B\Rightarrow \mu(A)\leq\mu(B)
\]

Для $c\in[0,1]$:
\[
\mu(-A+c)=\mu(A)
\]

\[
\mu(A+c)=\mu(A),\qquad x\in A
\]


\begin{note}
Мера на $[0,1]$: $\Delta\geq0$.
\end{note}

\[
E_k=E_0+r_k,
\]
где $r_k$ --- рациональное число из $[-1,1]$.

\[
[0,1]\subset \bigcup E_k\subset[-2,2]
\]

Множества $E_k$ попарно не пересекаются.

\[
1\leq \sum_{k=1}^{\infty}\Delta
\leq \mu([-2,2])
\]

Получаем противоречие.

\end{example}

\newpage

\section{Случайные величины и распределения}

\begin{definition}

Пусть $(\Omega,\mathcal{A},P)$ --- вероятностное пространство.

\textbf{Случайной величиной} называется отображение
\[
X:\Omega\to\mathbb{R},
\]
такое, что
\[
\forall B\in\mathcal{B}(\mathbb{R})
\qquad
X^{-1}(B)\in\mathcal{A}
\]

То есть $X$ является измеримым отображением.

\end{definition}

Пусть $X$ --- случайная величина на $(\Omega,\mathcal{A},P)$.

\begin{definition}

\textbf{Распределением случайной величины $X$} называется отображение
\[
P_X:\mathcal{B}(\mathbb{R})\to\mathbb{R},
\]
задаваемое равенством
\[
P_X(B)=P\bigl(X^{-1}(B)\bigr)
=P\bigl(\{\omega:X(\omega)\in B\}\bigr)
=P(X\in B)
\]

\end{definition}

\begin{note}

1) $P_X$ является вероятностной мерой. Таким образом,
\[
(\Omega,\mathcal{A},P)
\xrightarrow{\,X\,}
(\mathbb{R},\mathcal{B}(\mathbb{R}),P_X)
\]

2) Для случайных величин будем использовать обозначения
\[
X,Y,U,W,\ldots
\]
Также используются обозначения:
\[
\xi,\eta,\zeta,\ldots
\]

\end{note}

\begin{example}

Пусть
\[
\Omega=\{\text{орёл},\text{решка},\text{ребро}\},
\]
причём
\[
P(\text{орёл})=P(\text{решка})=\frac12,
\qquad
P(\text{ребро})=0
\]

Определим случайные величины $X$ и $Y$:
\[
X(\text{орёл})=0,\qquad
X(\text{решка})=1,\qquad
X(\text{ребро})=9982026,
\]
\[
Y(\text{решка})=0,\qquad
Y(\text{орёл})=1,\qquad
Y(\text{ребро})=9982026
\]

Тогда
\[
P(X=0)=\frac12=P(Y=1),
\]
\[
P(X=1)=\frac12=P(Y=0)
\]

Следовательно,
\[
P_X=P_Y,
\]
то есть
\[
X\overset{d}{=}Y
\]

При этом
\[
P(X=Y)=0,
\qquad
P(X\neq Y)=1,
\]
то есть $X$ и $Y$ почти наверное не равны.

\medskip

Рассмотрим модель геометрической вероятности
\[
([0,1],\mathcal{A},P),
\]
где
\[
\mathcal{A}
=
\left\{
A\subset[0,1]:
\operatorname{length}(A)\ \text{существует}
\right\},
\qquad
P(A)=\operatorname{length}(A)
\]

Пусть
\[
X(\omega)=1-\omega
\]

Тогда для $0\leq a\leq b\leq1$
\[
\begin{aligned}
P(a\leq X\leq b)
&=P(a\leq1-\omega\leq b)\\
&=P(1-b\leq\omega\leq1-a)\\
&=(1-a)-(1-b)\\
&=b-a
\end{aligned}
\]

\end{example}

\begin{definition}

Пусть $X$ --- случайная величина на $(\Omega,\mathcal{A},P)$, а $P_X$ --- её распределение.

\textbf{Функцией распределения} случайной величины $X$ называется функция
\[
F_X:\mathbb{R}\to\mathbb{R},
\]
задаваемая равенством
\[
F_X(t)
=
P_X((-\infty,t])
=
P(X\leq t)
\]

\end{definition}

\begin{theorem}{ - Свойства функции распределения}

Функция распределения $F$ обладает следующими свойствами:

\begin{enumerate}
    \item $F$ неубывает;
    \item $F$ непрерывна справа;
    \item
    $F(-\infty)=0,
    \qquad
    F(+\infty)=1$
    \item Для $a<b$
    \[
    P(a<X\leq b)=F_X(b)-F_X(a)
    \]

    \item Левый предел функции распределения удовлетворяет равенству
    \[
    F(t-0)=\lim_{s\uparrow t}F(s),
    \]
    причём
    \[
    P(X=t)=F(t)-F(t-0)
    \]

    \item Зная функцию распределения $F$, можно восстановить вероятности промежутков:
    \[
    P(a<X\leq b)=F(b)-F(a)
    \]
\end{enumerate}

\end{theorem}

\begin{proofbox}{монотонности}

Пусть $t_1<t_2$. Тогда
\[
(-\infty,t_1]\subset(-\infty,t_2]
\]

Следовательно,
\[
P_X((-\infty,t_1])
\leq
P_X((-\infty,t_2]),
\]
то есть
\[
F(t_1)\leq F(t_2)
\]

Таким образом, $F$ неубывает.

\end{proofbox}

\begin{proofbox}{предельных значений}

Докажем сначала, что
\[
F(+\infty)=1,
\qquad\text{то есть}\qquad
\lim_{t\to+\infty}F(t)=1
\]

Возьмём последовательность $\{a_n\}$ такую, что
\[
a_n\leq a_{n+1},
\qquad
a_n\to+\infty
\]

Положим
\[
A_n=(-\infty,a_n]
\]

Тогда
\[
A_n\subset A_{n+1},
\qquad
\bigcup_{n=1}^{\infty}A_n=\mathbb{R}
\]

По непрерывности вероятностной меры снизу
\[
F(a_n)=P_X(A_n)\longrightarrow P_X(\mathbb{R})=1
\]

Аналогично докажем, что
\[
F(-\infty)=0,
\qquad\text{то есть}\qquad
\lim_{t\to-\infty}F(t)=0
\]

Возьмём последовательность $\{b_n\}$ такую, что
\[
b_{n+1}\leq b_n,
\qquad
b_n\to-\infty
\]

Положим
\[
B_n=(-\infty,b_n]
\]

Тогда
\[
B_{n+1}\subset B_n,
\qquad
\bigcap_{n=1}^{\infty}B_n=\emptyset
\]

По непрерывности вероятностной меры сверху
\[
F(b_n)=P_X(B_n)\longrightarrow P_X(\emptyset)=0
\]

\end{proofbox}

\begin{proofbox}{непрерывности справа}

Необходимо показать, что
\[
F(t+0)=F(t)
\]

Пусть
\[
t_{n+1}\leq t_n,
\qquad
t_n\to0,
\qquad
t_n>0
\]

Тогда
\[
\begin{aligned}
F(t+t_n)-F(t)
&=P(X\leq t+t_n)-P(X\leq t)\\
&=P(t<X\leq t+t_n)
\end{aligned}
\]

Обозначим
\[
B_n=(t,t+t_n]
\]

Тогда
\[
B_{n+1}\subset B_n,
\qquad
\bigcap_{n=1}^{\infty}B_n=\emptyset
\]

По непрерывности вероятностной меры сверху
\[
P_X(B_n)\longrightarrow P_X(\emptyset)=0
\]

Следовательно,
\[
F(t+t_n)-F(t)\longrightarrow0,
\]
а значит,
\[
F(t+t_n)\longrightarrow F(t)
\]

Таким образом, $F$ непрерывна справа.

\end{proofbox}

\begin{proofbox}{свойства 4}

Так как
\[
\{X\leq b\}
=
\{X\leq a\}\sqcup\{a<X\leq b\},
\]
то
\[
F_X(b)
=
F_X(a)+P(a<X\leq b)
\]

Следовательно,
\[
P(a<X\leq b)=F_X(b)-F_X(a)
\]

\end{proofbox}

\begin{proofbox}{свойства 5}

Пусть $t_n>0$ и $t_n\to0$. Тогда
\[
F(t)-F(t-t_n)
=
P(t-t_n<X\leq t)
\]

При этом
\[
\bigcap_{n=1}^{\infty}(t-t_n,t]=\{t\}
\]

По непрерывности вероятностной меры сверху
\[
F(t)-F(t-0)=P(X=t)
\]

В частности,
\[
F(t-0)=F(t)
\quad\Longleftrightarrow\quad
P(X=t)=0
\]

\end{proofbox}

\begin{proofbox}{свойства 6}

Свойство непосредственно следует из первого:
\[
P(a<X\leq b)=F(b)-F(a)
\]

Следовательно, функция распределения определяет вероятности промежутков данного вида.

\end{proofbox}

\begin{theorem}{о достаточности функции распределения для описания вероятностного закона}

Пусть
\[
F:\mathbb{R}\to\mathbb{R}
\]
--- неубывающая функция, непрерывная справа и такая, что
\[
F(+\infty)=1,
\qquad
F(-\infty)=0
\]

Тогда существует вероятностное пространство $(\Omega,\mathcal{A},P)$ и случайная величина $X$ на $(\Omega,\mathcal{A},P)$ такие, что
\[
F_X=F
\]

\end{theorem}

\begin{proofbox}{существования случайной величины с заданной функцией распределения}

Рассмотрим алгебру множеств
\[
\mathcal{A}_0
=
\left\{
\bigsqcup_{k=1}^{m}(a_k,b_k]:
m<\infty
\right\}
\]

Определим на ней функцию
\[
\widetilde{P}
\left(
\bigsqcup_{k=1}^{m}(a_k,b_k]
\right)
=
\sum_{k=1}^{m}
\bigl(F(b_k)-F(a_k)\bigr)
\]

В частности,
\[
P(\ldots)\geq0,
\]
а также
\[
(-\infty,+\infty)\in\mathcal{A}_0
\]
и
\[
P((-\infty,+\infty))
=
F(+\infty)-F(-\infty)
=
1
\]

Далее необходимо доказать непрерывность построенной вероятностной меры для счётной аддитивности.

После этого применяем теорему о продолжении вероятностной меры и получаем вероятностную меру $\widetilde{P}$ на $\mathcal{B}(\mathbb{R})$, совпадающую с $P$ на $\mathcal{A}_0$.

Рассмотрим вероятностное пространство
\[
(\mathbb{R},\mathcal{B}(\mathbb{R}),P)
\]
и случайную величину
\[
X(\omega)=\omega
\]

Тогда по построению
\[
P(X\leq t)=P((-\infty,t])=F(t),
\]
то есть
\[
F_X=F
\]

\end{proofbox}

\newpage

\section{Типы случайных величин и распределений}

Пусть $(\Omega,\mathcal{A},P)$ --- вероятностное пространство, $X$ --- случайная величина, $P_X$ --- её распределение.

\begin{definition}

Случайная величина $X$ называется \textbf{дискретной}, если существует не более чем счётное множество $E\subset\mathbb{R}$ такое, что
\[
P(X\in E)=1
\]

\end{definition}

Пусть возможные значения $X$ равны $x_1,x_2,x_3,\ldots$, причём
\[
P(X=x_k)=p_k,
\qquad
\sum_k p_k=1
\]

Закон распределения можно записать в виде таблицы:
\[
\begin{array}{c|cccccc}
X & x_1 & x_2 & x_3 & \ldots & x_n & \ldots\\
\hline
P_X & p_1 & p_2 & p_3 & \ldots & p_n & \ldots
\end{array}
\]

При этом
\[
p(x)=
\begin{cases}
p_k, & x=x_k,\\
0, & \text{иначе}
\end{cases}
\]

Функция распределения дискретной случайной величины имеет вид
\[
F(t)=\sum_k p_k\,\mathds{1}(x_k\leq t)
\]

\begin{example}

\textbf{1. Биномиальное распределение}

Пусть
\[
X\sim \operatorname{Bin}(n,p),
\qquad
p\in(0,1),
\qquad
q=1-p,
\]
где $n$ --- количество испытаний, а $X$ --- количество успехов.

Тогда
\[
P(X=k)=C_n^k p^k q^{\,n-k},
\qquad
k=0,1,\ldots,n
\]

\medskip

\textbf{2. Распределение Бернулли}

Распределение Бернулли является частным случаем биномиального распределения:
\[
\operatorname{Bern}(p)=\operatorname{Bin}(1,p)
\]

\medskip

\textbf{3. Вырожденное распределение}

Пусть
\[
X\sim I(c)=I_c
\]

Тогда
\[
X\overset{\text{п.н.}}{=}c
\qquad\Longleftrightarrow\qquad
P(X=c)=1
\]

\medskip

\textbf{4. Геометрическое распределение}

Пусть $X$ --- количество неудач до первого успеха. Тогда
\[
P(X=k)=q^k p,
\qquad
k=0,1,2,\ldots
\]

Если же $X$ --- номер первого успеха, то
\[
P(X=k)=q^{k-1}p,
\qquad
k=1,2,\ldots
\]

\medskip

\textbf{5. Распределение Пуассона}

Пусть
\[
X\sim \operatorname{Pois}(\lambda),
\qquad
\lambda>0
\]

Тогда
\[
P(X=k)=e^{-\lambda}\frac{\lambda^k}{k!},
\qquad
k=0,1,2,\ldots
\]

\begin{theorem}{Пуассона}

Пусть
\[
X_n\sim\operatorname{Bin}(n,p_n),
\qquad
p_n=\frac{\lambda}{n}+o\left(\frac1n\right)
\]

Тогда для любого фиксированного $k$
\[
P(X_n=k)
=
C_n^k p_n^k(1-p_n)^{n-k}
\longrightarrow
e^{-\lambda}\frac{\lambda^k}{k!}
\]

\end{theorem}

\begin{proofbox}{предельного перехода к распределению Пуассона}

Имеем
\[
P(X_n=k)
=
C_n^k p_n^k(1-p_n)^{n-k}
\]

Так как
\[
p_n=\frac{\lambda}{n}+o\left(\frac1n\right),
\]
то при фиксированном $k$
\[
C_n^k p_n^k
=
\frac{n(n-1)\cdots(n-k+1)}{k!}\,p_n^k
\longrightarrow
\frac{\lambda^k}{k!}
\]

Кроме того,
\[
(1-p_n)^{n-k}
=
\frac{(1-p_n)^n}{(1-p_n)^k}
\longrightarrow e^{-\lambda}
\]

Следовательно,
\[
P(X_n=k)
\longrightarrow
e^{-\lambda}\frac{\lambda^k}{k!}
\]

\end{proofbox}

\textbf{6. Гипергеометрическое распределение}

Пусть
\[
X\sim HG(m,n,k),
\]
где $m$ --- общее количество деталей в корзине, $n$ --- количество ``хороших'' деталей, $k$ --- количество выбираемых без возвращения деталей.

Тогда
\[
P(X=j)
=
\frac{C_n^j C_{m-n}^{k-j}}{C_m^k}
\]

Если
\[
m\to\infty,
\qquad
n\to\infty,
\qquad
\frac{n}{m}\to p,
\]
то
\[
\frac{C_n^j C_{m-n}^{k-j}}{C_m^k}
\longrightarrow
C_k^j p^j q^{\,k-j},
\qquad
q=1-p,
\]
то есть в пределе получаем биномиальное распределение.

\medskip

\textbf{7. Отрицательное биномиальное распределение}

Пусть
\[
X\sim NB(r,p),
\]
где $X$ --- количество неудач до $r$-го успеха.

Тогда
\[
P(X=j)
=
C_{k+r-1}^{j}p^r q^j
\]

\medskip

\textbf{8. Дискретное равномерное распределение}

Пусть
\[
X\sim DU(x_1,\ldots,x_n)
\]

Тогда
\[
P(X=x_k)=\frac1n,
\qquad
k=1,\ldots,n
\]

\end{example}

\newpage

\section{Абсолютно непрерывные случайные величины и распределения}

\begin{definition}

Случайная величина $X$ и её распределение $P_X$ называются \textbf{абсолютно непрерывными}, если существует функция
\[
p:\mathbb{R}\to\mathbb{R},
\qquad
p(x)\geq0,
\]
такая, что для любого $B\in\mathcal{B}(\mathbb{R})$
\[
P(X\in B)=\int_B p(x)\,dx
\]

Функция $p(x)$ называется \textbf{плотностью распределения}, или PDF (от англ. \textit{probability density function}).

\end{definition}

\begin{props}{абсолютно непрерывного распределения}

1) Функция распределения имеет вид
\[
F(t)=P(X\leq t)
=
\int_{-\infty}^{t}p(x)\,dx
\]

В частности,
\[
F\in C(\mathbb{R}),
\]
и
\[
F'(t)=p(t)
\]
почти всюду.

2) Нормировка плотности:
\[
\int_{\mathbb{R}}p(x)\,dx=1
\]

3) Для любого $B\in\mathcal{B}(\mathbb{R})$
\[
P(X\in B)=\int_Bp(x)\,dx
\]

4) Для любого $c\in\mathbb{R}$
\[
P(X=c)
=
\int_{\{c\}}p(x)\,dx
=
0
\]

5) При $h\to0$
\[
F(x_0+h)-F(x_0)
=
p(x_0)h+o(h),
\]
если $p$ непрерывна в точке $x_0$. Поэтому
\[
P(x_0\leq X\leq x_0+h)
\approx
p(x_0)h
\]

\end{props}

Основные аналогии:

\[
\begin{array}{c|c|c}
 & \text{Дискретное} & \text{Абсолютно непрерывное}\\
\hline
\text{Описание}
&
P(X=x_i)
&
\text{Плотность }p(x)
\\[2mm]
\text{Нормировка}
&
\displaystyle\sum_k p(x_k)=1
&
\displaystyle\int_{\mathbb{R}}p(x)\,dx=1
\\[4mm]
P(a<X<b)
&
\displaystyle\sum_{a<x_k<b}p(x_k)
&
\displaystyle\int_a^b p(x)\,dx
\end{array}
\]

\subsection{Равномерное распределение}

\begin{definition}

Случайная величина $X$ имеет \textbf{равномерное распределение} на $[a,b]$,
\[
X\sim U[a,b],
\]
если её плотность равна
\[
p_X(t)
=
\frac{1}{b-a}\mathds{1}(a\leq t\leq b)
\]

\end{definition}

Функция распределения имеет вид
\[
F_X(t)
=
\begin{cases}
0, & t<a,\\[1mm]
\dfrac{t-a}{b-a}, & a\leq t\leq b,\\[2mm]
1, & t>b
\end{cases}
\]

Пусть
\[
Y=cX+d,
\qquad
c>0
\]

Тогда
\[
Y\sim U[ca+d,cb+d]
\]

Действительно,
\[
F_Y(t)
=
P(cX+d\leq t)
=
P\left(X\leq\frac{t-d}{c}\right)
=
F_X\left(\frac{t-d}{c}\right)
\]

Следовательно,
\[
F_Y(t)
=
\begin{cases}
0, & \dfrac{t-d}{c}<a,\\[2mm]
\dfrac{\frac{t-d}{c}-a}{b-a},
& a\leq\dfrac{t-d}{c}\leq b,\\[3mm]
1, & \dfrac{t-d}{c}>b
\end{cases}
\]

Эквивалентно,
\[
F_Y(t)
=
\begin{cases}
0, & t<ca+d,\\[1mm]
\dfrac{t-(ca+d)}{c(b-a)},
& ca+d\leq t\leq cb+d,\\[3mm]
1, & t>cb+d
\end{cases}
\]

\subsection{Нормальное распределение}

\begin{definition}

Случайная величина $X$ имеет \textbf{нормальное распределение} с параметрами $\mu\in\mathbb{R}$ и $\sigma>0$,
\[
X\sim\mathcal{N}(\mu,\sigma^2),
\]
если её плотность имеет вид
\[
p(x)
=
\frac{1}{\sqrt{2\pi\sigma^2}}
\exp\left(
-\frac{(x-\mu)^2}{2\sigma^2}
\right)
\]

Параметр $\mu$ определяет положение распределения, а $\sigma^2$ --- его масштаб.

\end{definition}

Распределение
\[
\mathcal{N}(0,1)
\]
называется \textbf{стандартным нормальным распределением}.

\begin{note}

Для нормального распределения $\mu$
является математическим ожиданием, а $\sigma^2$ --- дисперсией.

\end{note}

Пусть
\[
X\sim\mathcal{N}(\mu,\sigma^2),
\qquad
Y=aX+b,
\qquad
a\neq0
\]

Тогда
\[
Y\sim\mathcal{N}(a\mu+b,a^2\sigma^2)
\]

Действительно,
\[
F_Y(t)
=
P(aX+b\leq t)
\]

При $a>0$
\[
F_Y(t)
=
P\left(X\leq\frac{t-b}{a}\right)
=
F_X\left(\frac{t-b}{a}\right),
\]
поэтому
\[
\begin{aligned}
p_Y(t)
&=
\frac1a
p_X\left(\frac{t-b}{a}\right)\\
&=
\frac{1}{\sqrt{2\pi\sigma^2a^2}}
\exp\left(
-\frac{(t-b-a\mu)^2}{2a^2\sigma^2}
\right)
\end{aligned}
\]

Следовательно,
\[
Y\sim\mathcal{N}(a\mu+b,a^2\sigma^2)
\]

\begin{note}

Стандартизация нормальной случайной величины:
\[
\frac{X-\mu}{\sigma}\sim\mathcal{N}(0,1)
\]

Если
\[
Y\sim\mathcal{N}(0,1),
\]
то
\[
\sigma Y+\mu\sim\mathcal{N}(\mu,\sigma^2)
\]

\end{note}

\subsection{Принцип трёх сигм}

Пусть
\[
X\sim\mathcal{N}(\mu,\sigma^2)
\]

Тогда
\[
\begin{aligned}
P(\mu-c\sigma\leq X\leq\mu+c\sigma)
&=
P\left(
-c\leq\frac{X-\mu}{\sigma}\leq c
\right)\\
&=
\Phi(c)-\Phi(-c)\\
&=
2\Phi(c)-1,
\end{aligned}
\]
где $\Phi$ --- функция распределения стандартного нормального закона.

Используется симметрия
\[
\Phi(-x)=1-\Phi(x),
\]
где
\[
\Phi(x)
=
\frac{1}{\sqrt{2\pi}}
\int_{-\infty}^{x}e^{-t^2/2}\,dt
\]

\begin{note}

Значения функции $\Phi$ обычно берутся из таблицы числового закона.

\end{note}

Обозначим
\[
\Phi_0(t)
=
\frac{1}{\sqrt{2\pi}}
\int_0^t e^{-x^2/2}\,dx
\]

Это \textbf{функция Лапласа}. Для неё
\[
\Phi_0(-t)=-\Phi_0(t),
\]
а функция распределения стандартного нормального закона выражается как
\[
\Phi(t)=\frac12+\Phi_0(t)
\]

\textbf{Функция ошибок} определяется равенством
\[
\operatorname{erf}(x)
=
\frac{2}{\sqrt{\pi}}
\int_0^x e^{-t^2}\,dt
\]

После замены
\[
t=\frac{u}{\sqrt2},
\qquad
dt=\frac{du}{\sqrt2},
\]
получаем
\[
\operatorname{erf}(x)
=
2\Phi_0(\sqrt2\,x)
\]

Также нормальную случайную величину можно записать в виде
\[
Y=\mu+\xi,
\qquad
\xi\sim\mathcal{N}(0,\sigma^2)
\]

\subsection{Экспоненциальное распределение}

\begin{definition}

Случайная величина $X$ имеет \textbf{экспоненциальное распределение} с параметром $\lambda>0$,
\[
X\sim\operatorname{Exp}(\lambda),
\]
если её плотность равна
\[
p(x)
=
\lambda e^{-\lambda x}\mathds{1}(x\geq0)
\]

Функция распределения:
\[
F(x)
=
(1-e^{-\lambda x})\mathds{1}(x\geq0)
\]

\end{definition}

\begin{note}

В другой параметризации экспоненциального распределения записывают
\[
p(x)
=
\frac1\lambda e^{-x/\lambda}\mathds{1}(x\geq0),
\]
и тогда
\[
F(x)
=
(1-e^{-x/\lambda})\mathds{1}(x\geq0)
\]

\end{note}

\begin{example}

Пусть моменты прихода запросов обозначены через
\[
\tau_1,\tau_2,\tau_3,\ldots
\]

Тогда интервалы между последовательными моментами прихода изображаются как
\[
[0,\tau_1],\quad
[\tau_1,\tau_2],\quad
[\tau_2,\tau_3],\ldots
\]

Если $T_k$ --- время между последовательными запросами, то в рассматриваемой модели
\[
T_k\sim\operatorname{Exp}(\lambda)
\]

\end{example}

\subsection{Гамма-распределение}

\begin{definition}

Случайная величина $X$ имеет \textbf{гамма-распределение} с параметрами $\alpha,\lambda>0$,
\[
X\sim\Gamma(\alpha,\lambda),
\]
если её плотность равна
\[
p(x)
=
\frac{\lambda^\alpha x^{\alpha-1}e^{-\lambda x}}
{\Gamma(\alpha)}
\mathds{1}(x\geq0)
\]

При $\alpha=1$
\[
\Gamma(1,\lambda)=\operatorname{Exp}(\lambda)
\]

\end{definition}

\subsection{Распределение Коши}

\begin{definition}

Случайная величина $X$ имеет \textbf{распределение Коши} с параметрами $\mu$ и $\gamma>0$, если её плотность имеет вид
\[
p(x)
=
\frac{1}{\pi\gamma}
\frac{1}{
1+\left(\dfrac{x-\mu}{\gamma}\right)^2
}
\]

\end{definition}

Функция распределения:

\[
F_X(t)
=
P(X\leq t)
=
\int_{-\infty}^{t}p(x)\,dx
\]

Сделаем замену
\[
u=\frac{x-\mu}{\gamma},
\qquad
dx=\gamma\,du
\]

Получаем
\[
\begin{aligned}
F_X(t)
&=
\int_{-\infty}^{t}
\frac{1}{\pi\gamma}
\frac{1}{
1+\left(\dfrac{x-\mu}{\gamma}\right)^2
}\,dx\\
&=
\frac{1}{\pi}
\int_{-\infty}^{\frac{t-\mu}{\gamma}}
\frac{du}{1+u^2}\\
&=
\frac{1}{\pi}
\left.
\arctan u
\right|_{-\infty}^{\frac{t-\mu}{\gamma}}\\
&=
\frac{1}{\pi}
\left(
\arctan\frac{t-\mu}{\gamma}
+\frac{\pi}{2}
\right)
\end{aligned}
\]

Следовательно,
\[
\boxed{
F_X(t)
=
\frac12+
\frac{1}{\pi}
\arctan\left(\frac{t-\mu}{\gamma}\right)
}
\]

\subsection{Плотность преобразованной случайной величины}

\begin{theorem}{о плотности преобразования}

Пусть $X$ --- абсолютно непрерывная случайная величина с плотностью $p_X$, и
\[
Y=g(X),
\]
где $g$ --- гладкая биекция.

Если $g$ строго монотонна, то
\[
p_Y(t)
=
p_X\bigl(g^{-1}(t)\bigr)
\left|
\left(g^{-1}\right)'(t)
\right|
\]

В частности,
\[
p_Y(t)
=
p_X\bigl(g^{-1}(t)\bigr)
\frac{1}{\left|g'(t)\right|}
\]

\end{theorem}

\begin{proofbox}{плотности преобразованной случайной величины}

Пусть сначала $g$ строго возрастает. Тогда
\[
\begin{aligned}
F_Y(t)
&=
P(g(X)\leq t)\\
&=
P(X\leq g^{-1}(t))\\
&=
F_X(g^{-1}(t))
\end{aligned}
\]

Дифференцируя, получаем
\[
p_Y(t)
=
p_X(g^{-1}(t))
\left(g^{-1}\right)'(t)
\]

Для строго убывающей функции аналогично возникает модуль производной, поэтому в общем случае
\[
p_Y(t)
=
p_X(g^{-1}(t))
\left|
\left(g^{-1}\right)'(t)
\right|
\]

\end{proofbox}

Если
\[
g: \bigsqcup_x([a_x, b_x])\longrightarrow \mathbb{R}
\]
и на каждом интервале монотонности g гладкая биекция, то
\[
p_Y(t)
=
\sum_{k=1}^{n(t)}
p_X\bigl(g_k^{-1}(t)\bigr)
\left|
\left(g_k^{-1}\right)'(t)
\right|,
\]
где
\[
n(t)=\#\{x:g(x)=t\}
\]

\newpage

\section{Случайные векторы. Совместные распределения}

\begin{definition}

\textbf{Носителем распределения} случайной величины $X$ называется замкнутое множество $E$ такое, что
\[
P(X\in E)=1
\]
и для любого собственного замкнутого подмножества $E'\subsetneq E$
\[
P(X\in E')<1
\]

\end{definition}

Для дискретного распределения
\[
E=\{x:p_X(x)>0\}
\]

Для абсолютно непрерывного распределения
\[
E=\operatorname{cl}\{x:p_X(x)>0\}
\]

\begin{definition}

Пусть
\[
X=(X_1,\ldots,X_n)
\]
где $X_1,\ldots,X_n$ --- случайные величины. Тогда $X$ называется \textbf{случайным вектором}

\end{definition}

В дискретном двумерном случае совместное распределение задаётся вероятностями
\[
p_{ij}=P(X=x_i,Y=y_j)
\]

Маргинальные распределения имеют вид
\[
P(X=x_i)=\sum_j p_{ij},
\qquad
P(Y=y_j)=\sum_i p_{ij}
\]

Если
\[
U=g(X,Y)
\]
то
\[
P(U=u_k)
=
\sum_{\substack{i,j\\g(x_i,y_j)=u_k}}p_{ij}
\]

\subsection{Абсолютно непрерывный случайный вектор}

\begin{definition}

Случайный вектор
\[
X=(X_1,\ldots,X_n)
\]
называется \textbf{абсолютно непрерывным}, если существует функция
\[
p:\mathbb{R}^n\to\mathbb{R}
\]
такая, что для любого $B\in\mathcal{B}(\mathbb{R}^n)$
\[
P(X\in B)
=
\int_B p(x_1,\ldots,x_n)\,dx_1\cdots dx_n
\]

Функция $p$ называется \textbf{совместной плотностью} случайного вектора

\end{definition}

Для двумерного случайного вектора $(X,Y)$ маргинальные плотности равны
\[
p_X(x)=\int_{\mathbb{R}}p(x,y)\,dy
\]
и
\[
p_Y(y)=\int_{\mathbb{R}}p(x,y)\,dx
\]

Если
\[
U=g(X,Y)
\]
то для любого $B\in\mathcal{B}(\mathbb{R})$
\[
P(U\in B)
=
\iint_{\{(x,y):g(x,y)\in B\}}p(x,y)\,dx\,dy
\]

\begin{example}

Пусть совместная плотность имеет вид
\[
p(x,y)
=
c e^{-x}e^{-2y}
\mathds{1}\left(x\geq0,\ y\geq0,\ y\leq\frac1x\right)
\]

Найдём нормировочную константу $c$. Из условия
\[
\iint_{\mathbb{R}^2}p(x,y)\,dx\,dy=1
\]
получаем
\[
\begin{aligned}
1
&=
c\int_0^{+\infty}e^{-x}
\int_0^{1/x}e^{-2y}\,dy\,dx\\
&=
c\int_0^{+\infty}
e^{-x}\frac{1-e^{-2/x}}{2}\,dx
\end{aligned}
\]

Следовательно,
\[
c^{-1}
=
\frac12\int_0^{+\infty}
e^{-x}\left(1-e^{-2/x}\right)\,dx
\]
и численно
\[
c\approx 2.14
\]

Маргинальная плотность $X$:
\[
\begin{aligned}
p_X(x)
&=
\int_{\mathbb{R}}
c e^{-x}e^{-2y}
\mathds{1}\left(y\geq0,\ x\geq0,\ y\leq\frac1x\right)\,dy\\
&=
\frac c2 e^{-x}\left(1-e^{-2/x}\right)\mathds{1}(x>0)
\end{aligned}
\]

Маргинальная плотность $Y$:
\[
\begin{aligned}
p_Y(y)
&=
\int_{\mathbb{R}}
c e^{-x}e^{-2y}
\mathds{1}\left(y\geq0,\ x\geq0,\ y\leq\frac1x\right)\,dx\\
&=
c e^{-2y}\int_0^{1/y}e^{-x}\,dx\,
\mathds{1}(y>0)\\
&=
c e^{-2y}\left(1-e^{-1/y}\right)\mathds{1}(y>0)
\end{aligned}
\]

\medskip

\textbf{Равномерное распределение.}

Пусть $D\subset\mathbb{R}^n$. Равномерное распределение на $D$ имеет плотность
\[
p(x)
=
\frac{1}{\operatorname{Vol}(D)}
\mathds{1}(x\in D)
\]

\medskip

\textbf{Многомерное нормальное распределение.}

Стандартный гауссовский случайный вектор
\[
X=(X_1,\ldots,X_n)
\]
имеет совместную плотность
\[
p(x_1,\ldots,x_n)
=
\frac{1}{(2\pi)^{n/2}}
\exp\left(
-\frac12\sum_{i=1}^n x_i^2
\right)
\]

То есть его компоненты являются стандартными нормальными случайными величинами

Пусть
\[
X\sim\mathcal{N}(0,I_n),
\qquad
Y=AX+b
\]
где
\[
A\in M_{m\times n},
\qquad
b\in\mathbb{R}^m
\]

Тогда
\[
Y\sim\mathcal{N}(b,AA^T)
\]

В общем случае многомерное нормальное распределение задаётся параметрами
\[
\mu\in\mathbb{R}^n,
\qquad
\Sigma\in M_{n\times n},
\qquad
\Sigma=\Sigma^T,
\qquad
\Sigma\geq0
\]

Если
\[
\Sigma=U\Lambda U^T,
\qquad
\Lambda=\operatorname{diag}(\lambda_1,\ldots,\lambda_n),
\qquad
\lambda_i\geq0
\]
то можно положить
\[
\sqrt{\Sigma}
=
U\sqrt{\Lambda}\,U^T
\]
где
\[
\sqrt{\Lambda}
=
\operatorname{diag}(\sqrt{\lambda_1},\ldots,\sqrt{\lambda_n})
\]

При
\[
A=\sqrt{\Sigma},
\qquad
b=\mu
\]
получаем случайный вектор с распределением
\[
Y\sim\mathcal{N}(\mu,\Sigma)
\]

\end{example}

\begin{note}

Если $\Sigma>0$, то многомерное нормальное распределение имеет плотность
\[
p(x_1,\ldots,x_n)
=
\frac{1}{(2\pi)^{n/2}\sqrt{\det\Sigma}}
\exp\left(
-\frac12(x-\mu)^T\Sigma^{-1}(x-\mu)
\right)
\]

\end{note}

\subsection{Функция распределения случайного вектора}

\begin{definition}

\textbf{Функцией распределения} случайного вектора
\[
X=(X_1,\ldots,X_n)
\]
называется функция
\[
F(t_1,\ldots,t_n)
=
P(X_1\leq t_1,\ldots,X_n\leq t_n)
\]

\end{definition}

\begin{props}{функции распределения случайного вектора}

\begin{enumerate}
    \item
    \[
    \lim_{\substack{t_1\to-\infty\\\text{или хотя бы одного }t_i\to-\infty}}
    F(t_1,\ldots,t_n)=0
    \]
    и
    \[
    \lim_{\substack{t_1\to+\infty\\\cdots\\t_n\to+\infty}}
    F(t_1,\ldots,t_n)=1
    \]

    \item $F$ непрерывна справа по каждой переменной

    \item Для любого прямоугольника
    \[
    (a_1,b_1]\times\cdots\times(a_n,b_n]
    \]
    соответствующее конечное приращение функции распределения неотрицательно
    \[
    \Delta_{a_1,b_1}\cdots\Delta_{a_n,b_n}F\geq0
    \]
\end{enumerate}

\end{props}

Если случайный вектор абсолютно непрерывен, то
\[
F(t_1,\ldots,t_n)
=
\int_{-\infty}^{t_1}\cdots
\int_{-\infty}^{t_n}
p(x_1,\ldots,x_n)\,dx_n\cdots dx_1
\]
и
\[
p(x_1,\ldots,x_n)
=
\frac{\partial^n F(x_1,\ldots,x_n)}
{\partial x_1\cdots\partial x_n}
\]

\newpage

\section{Независимые случайные величины}

\begin{definition}

Случайные величины $X_1,\ldots,X_n$ называются \textbf{независимыми}, если для любых
\[
B_1,\ldots,B_n\in\mathcal{B}(\mathbb{R})
\]
выполняется
\[
P(X_1\in B_1,\ldots,X_n\in B_n)
=
\prod_{j=1}^n P(X_j\in B_j)
\]

Последовательность $(X_k)_{k=1}^{\infty}$ называется независимой, если для любого $n$ случайные величины
\[
X_1,\ldots,X_n
\]
независимы

\end{definition}

\begin{theorem}{Критерий независимости через функции распределения}

Случайные величины $X_1,\ldots,X_n$ независимы тогда и только тогда, когда
\[
F(t_1,\ldots,t_n)
=
\prod_{j=1}^n F_j(t_j)
\]

\end{theorem}

\begin{proofbox}{критерия независимости через функции распределения}

Если $X_1,\ldots,X_n$ независимы, то, взяв
\[
B_j=(-\infty,t_j]
\]
получаем
\[
\begin{aligned}
F(t_1,\ldots,t_n)
&=
P(X_1\leq t_1,\ldots,X_n\leq t_n)\\
&=
\prod_{j=1}^nP(X_j\leq t_j)\\
&=
\prod_{j=1}^nF_j(t_j)
\end{aligned}
\]

Обратное утверждение следует из того, что множества вида $(-\infty,t]$ порождают борелевскую $\sigma$-алгебру

\end{proofbox}

\begin{note}

Идея критерия для двумерного случая состоит в рассмотрении малых прямоугольников. Например,
\[
P(a_1<X_1\leq a_1+\sigma_1,\,
  a_2<X_2\leq a_2+\sigma_2)
\]
можно записать как
\[
\Delta_{a_1,a_1+\sigma_1}
\Delta_{a_2,a_2+\sigma_2}
F(a_1,a_2)
\]

Если
\[
F(t_1,t_2)=F_{X_1}(t_1)F_{X_2}(t_2)
\]
то
\[
\begin{aligned}
&P(a_1<X_1\leq a_1+\sigma_1,\,
a_2<X_2\leq a_2+\sigma_2)\\
&=
\left[
F_{X_1}(a_1+\sigma_1)-F_{X_1}(a_1)
\right]
\left[
F_{X_2}(a_2+\sigma_2)-F_{X_2}(a_2)
\right]
\end{aligned}
\]

\end{note}

\begin{theorem}{Критерий независимости дискретных случайных величин}

Пусть $X_1,\ldots,X_n$ --- дискретные случайные величины. Тогда они независимы тогда и только тогда, когда для любых возможных значений $x_{1,j_1},\ldots,x_{n,j_n}$
\[
P(X_1=x_{1,j_1},\ldots,X_n=x_{n,j_n})
=
\prod_{i=1}^nP(X_i=x_{i,j_i})
\]

\end{theorem}

\begin{proofbox}{критерия независимости дискретных случайных величин}

Необходимость непосредственно следует из определения независимости при выборе одноточечных борелевских множеств
\[
B_i=\{x_{i,j_i}\}
\]

Обратно, любое событие вида $\{X_i\in B_i\}$ для дискретной случайной величины является объединением соответствующих событий $\{X_i=x_{i,j}\}$. Суммируя равенство из условия по всем возможным значениям, принадлежащим $B_1,\ldots,B_n$, получаем
\[
P(X_1\in B_1,\ldots,X_n\in B_n)
=
\prod_{i=1}^nP(X_i\in B_i)
\]

Следовательно, $X_1,\ldots,X_n$ независимы

\end{proofbox}

\subsection{Совместное распределение независимых случайных величин}

Если $X_1,\ldots,X_n$ независимы и имеют плотности $p_{X_1},\ldots,p_{X_n}$, то их совместная плотность равна
\[
p(x_1,\ldots,x_n)
=
\prod_{j=1}^n p_{X_j}(x_j)
\]

В частности, стандартный гауссовский вектор имеет плотность
\[
\begin{aligned}
p(x_1,\ldots,x_n)
&=
\frac{1}{(2\pi)^{n/2}}
\exp\left(-\frac12x^Tx\right)\\
&=
\prod_{j=1}^n
\frac{1}{\sqrt{2\pi}}
\exp\left(-\frac{x_j^2}{2}\right)
\end{aligned}
\]
поэтому его компоненты независимы и имеют распределение $\mathcal{N}(0,1)$

Если
\[
X\sim\mathcal{N}
\left(
\mu,
\operatorname{diag}(\sigma_1^2,\ldots,\sigma_n^2)
\right)
\]
то
\[
p(x_1,\ldots,x_n)
=
\frac{1}{\sqrt{(2\pi)^n\sigma_1^2\cdots\sigma_n^2}}
\exp\left(
-\frac12
(x-\mu)^T
\begin{pmatrix}
\sigma_1^{-2} & & 0\\
& \ddots &\\
0 & & \sigma_n^{-2}
\end{pmatrix}
(x-\mu)
\right)
\]

\subsection{Распределение суммы независимых случайных величин}

Пусть $X$ и $Y$ независимы.

В дискретном случае
\[
P(X+Y=i)
=
\sum_jP(X=i-j)P(Y=j)
\]

Для абсолютно непрерывных случайных величин
\[
\begin{aligned}
F_{X+Y}(t)
&=
P(X+Y\leq t)\\
&=
\iint_{x+y\leq t}
p_X(x)p_Y(y)\,dx\,dy\\
&=
\int_{-\infty}^{+\infty}
p_Y(y)
\left(
\int_{-\infty}^{t-y}p_X(x)\,dx
\right)dy\\
&=
\int_{-\infty}^{+\infty}
p_Y(y)F_X(t-y)\,dy
\end{aligned}
\]

Дифференцируя по $t$, получаем
\[
p_{X+Y}(t)
=
\int_{-\infty}^{+\infty}
p_Y(y)p_X(t-y)\,dy
\]

То есть плотность суммы независимых случайных величин является свёрткой их плотностей
\[
p_{X+Y}=p_X*p_Y
\]

\newpage

\section{Числовые характеристики случайных величин}

\subsection{Одномерные}

\begin{definition}

Математическое ожидание дискретной случайной величины $X$ задаётся формулой
\[
EX
=
\sum_k x_k P(X=x_k)
\]

Если $X$ имеет плотность $p_X$, то
\[
EX
=
\int_{\mathbb{R}}x p_X(x)\,dx
\]

\end{definition}

\begin{props}{математического ожидания}

1) Линейность:
\[
E(aX+bY)
=
aEX+bEY
\]

\medskip

2) Если $X$ и $Y$ независимы, то
\[
E(XY)
=
EX\,EY
\]

\medskip

3) Для неотрицательной случайной величины $X$ и $c>0$ выполняется неравенство Маркова:
\[
P(X\geq c)
\leq
\frac{EX}{c}
\]

\medskip

4) Для неотрицательной случайной величины $X$:
\[
EX
=
\int_0^{+\infty}P(X\geq r)\,dr
\]

\end{props}

\begin{proofbox}{свойства 1}

Для дискретных случайных величин
\[
\begin{aligned}
E(aX+bY)
&=
\sum_{x,y}(ax+by)P(X=x,Y=y)\\
&=
a\sum_{x,y}xP(X=x,Y=y)
+
b\sum_{x,y}yP(X=x,Y=y)\\
&=
aEX+bEY
\end{aligned}
\]

Для абсолютно непрерывных случайных величин доказательство аналогично следует из линейности интеграла:
\[
\begin{aligned}
E(aX+bY)
&=
\iint_{\mathbb{R}^2}
(ax+by)p_{X,Y}(x,y)\,dx\,dy\\
&=
aEX+bEY
\end{aligned}
\]

\end{proofbox}

\begin{proofbox}{свойства 2}

Пусть $X$ и $Y$ независимы. Тогда их совместное распределение является произведением маргинальных распределений

В абсолютно непрерывном случае
\[
p_{X,Y}(x,y)=p_X(x)p_Y(y)
\]

Поэтому
\[
\begin{aligned}
E(XY)
&=
\iint_{\mathbb{R}^2}xy\,p_X(x)p_Y(y)\,dx\,dy\\
&=
\left(\int_{\mathbb{R}}x p_X(x)\,dx\right)
\left(\int_{\mathbb{R}}y p_Y(y)\,dy\right)\\
&=
EX\,EY
\end{aligned}
\]

В дискретном случае рассуждение аналогично:
\[
\begin{aligned}
E(XY)
&=
\sum_{x,y}xyP(X=x)P(Y=y)\\
&=
EX\,EY
\end{aligned}
\]

\end{proofbox}

\begin{proofbox}{свойства 3}

Так как $X\geq0$, то
\[
X\geq c\,\mathds{1}(X\geq c)
\]

Следовательно,
\[
EX
\geq
cE\mathds{1}(X\geq c)
=
cP(X\geq c)
\]

Отсюда
\[
P(X\geq c)
\leq
\frac{EX}{c}
\]

\end{proofbox}

\begin{proofbox}{свойства 4}

Для неотрицательной случайной величины
\[
X(\omega)
=
\int_0^{+\infty}
\mathds{1}(X(\omega)\geq r)\,dr
\]

Интегрируя по $\omega$ и меняя порядок интегрирования, получаем
\[
\begin{aligned}
EX
&=
E\left[
\int_0^{+\infty}
\mathds{1}(X\geq r)\,dr
\right]\\
&=
\int_0^{+\infty}
E\mathds{1}(X\geq r)\,dr\\
&=
\int_0^{+\infty}
P(X\geq r)\,dr
\end{aligned}
\]

\end{proofbox}

\begin{definition}

\textbf{Дисперсией} случайной величины $X$ называется
\[
\operatorname{Var}X
=
E(X-EX)^2
=
EX^2-(EX)^2
\]

Если $X$ имеет плотность $p_X$, то
\[
\operatorname{Var}X
=
\int_{\mathbb{R}}x^2p_X(x)\,dx
-
\left(
\int_{\mathbb{R}}xp_X(x)\,dx
\right)^2
\]

\end{definition}

\begin{props}{дисперсии}

1) Дисперсия равна нулю тогда и только тогда, когда случайная величина почти наверное постоянна:
\[
\operatorname{Var}X=0
\quad\Longleftrightarrow\quad
X=c\quad\text{почти наверное}
\]

\medskip

2) Для $a,b\in\mathbb{R}$
\[
\operatorname{Var}(aX+b)
=
a^2\operatorname{Var}X
\]

\medskip

3) Для случайных величин $X_1,\ldots,X_n$ и чисел $c_1,\ldots,c_n$
\[
\operatorname{Var}\left(\sum_{i=1}^n c_iX_i\right)
=
\sum_{i=1}^n c_i^2\operatorname{Var}X_i
+
2\sum_{i<j}c_ic_j\operatorname{Cov}(X_i,X_j)
\]

\medskip

4) Неравенство Чебышёва:
\[
P\left(|X-EX|\geq\varepsilon\right)
\leq
\frac{\operatorname{Var}X}{\varepsilon^2}
\]

\medskip

5) Математическое ожидание минимизирует средний квадрат отклонения:
\[
EX
=
\operatorname*{arg\,min}_{c\in\mathbb{R}}
E(X-c)^2
\]

\end{props}

\begin{proofbox}{свойства 1}

По определению
\[
\operatorname{Var}X
=
E(X-EX)^2
\]

Так как $(X-EX)^2\geq0$, то
\[
E(X-EX)^2=0
\]
тогда и только тогда, когда
\[
(X-EX)^2=0
\quad\text{почти наверное}
\]

Следовательно,
\[
X=EX
\quad\text{почти наверное}
\]

Обратно, если $X=c$ почти наверное, то
\[
EX=c
\]
и
\[
\operatorname{Var}X
=
E(X-c)^2
=
0
\]

\end{proofbox}

\begin{proofbox}{свойства 2}

Имеем
\[
E(aX+b)=aEX+b
\]

Поэтому
\[
\begin{aligned}
\operatorname{Var}(aX+b)
&=
E\left(aX+b-aEX-b\right)^2\\
&=
E\left(a(X-EX)\right)^2\\
&=
a^2\operatorname{Var}X
\end{aligned}
\]

\end{proofbox}

\begin{proofbox}{свойства 3}

Обозначим
\[
Z=\sum_{i=1}^n c_iX_i
\]

Тогда
\[
Z-EZ
=
\sum_{i=1}^n c_i(X_i-EX_i)
\]

Следовательно,
\[
\begin{aligned}
\operatorname{Var}Z
&=
E\left(
\sum_{i=1}^n c_i(X_i-EX_i)
\right)^2\\
&=
\sum_{i=1}^n
c_i^2E(X_i-EX_i)^2\\
&\quad+
2\sum_{i<j}
c_ic_j
E\left[(X_i-EX_i)(X_j-EX_j)\right]\\
&=
\sum_{i=1}^n c_i^2\operatorname{Var}X_i
+
2\sum_{i<j}c_ic_j\operatorname{Cov}(X_i,X_j)
\end{aligned}
\]

\end{proofbox}

\begin{proofbox}{свойства 4}

Применим неравенство Маркова к неотрицательной случайной величине
\[
Y=(X-EX)^2
\]

Тогда
\[
\begin{aligned}
P(|X-EX|\geq\varepsilon)
&=
P\left((X-EX)^2\geq\varepsilon^2\right)\\
&\leq
\frac{E(X-EX)^2}{\varepsilon^2}\\
&=
\frac{\operatorname{Var}X}{\varepsilon^2}
\end{aligned}
\]

\end{proofbox}

\begin{proofbox}{свойства 5}

Для любого $c\in\mathbb{R}$
\[
X-c
=
(X-EX)+(EX-c)
\]

Следовательно,
\[
\begin{aligned}
E(X-c)^2
&=
E(X-EX)^2
+
2(EX-c)E(X-EX)
+
(EX-c)^2\\
&=
\operatorname{Var}X+(EX-c)^2
\end{aligned}
\]

Так как
\[
(EX-c)^2\geq0
\]
минимум достигается при
\[
c=EX
\]

\end{proofbox}

\begin{theorem}{о числовых характеристиках стандартных распределений}

1) Если $X\sim U[a,b]$, то
\[
EX=\frac{a+b}{2},
\qquad
\operatorname{Var}X=\frac{(b-a)^2}{12}
\]

\medskip

2) Если $X\sim\mathcal{N}(\mu,\sigma^2)$, то
\[
EX=\mu,
\qquad
\operatorname{Var}X=\sigma^2
\]

\medskip

3) Если $X\sim\operatorname{Exp}(\lambda)$ и
\[
p_X(x)=\lambda e^{-\lambda x}\mathds{1}(x\geq0)
\]
то
\[
EX=\frac{1}{\lambda},
\qquad
\operatorname{Var}X=\frac{1}{\lambda^2}
\]

\medskip

4) Если $X\sim\operatorname{Gamma}(\alpha,\lambda)$ и
\[
p_X(x)
=
\frac{\lambda^\alpha x^{\alpha-1}}{\Gamma(\alpha)}
e^{-\lambda x}\mathds{1}(x\geq0)
\]
то
\[
EX=\frac{\alpha}{\lambda},
\qquad
\operatorname{Var}X=\frac{\alpha}{\lambda^2}
\]

\medskip

5) Для стандартного распределения Коши математическое ожидание не существует, поскольку интеграл
\[
\int_{-\infty}^{+\infty}
\frac{x}{\pi(1+x^2)}\,dx
\]
не сходится абсолютно

\end{theorem}

\begin{proofbox}{свойства 1}

Для $X\sim U[a,b]$
\[
p_X(x)
=
\frac{1}{b-a}\mathds{1}(a\leq x\leq b)
\]

Поэтому
\[
EX
=
\frac{1}{b-a}\int_a^b x\,dx
=
\frac{a+b}{2}
\]

Кроме того,
\[
EX^2
=
\frac{1}{b-a}\int_a^b x^2\,dx
=
\frac{a^2+ab+b^2}{3}
\]

Следовательно,
\[
\operatorname{Var}X
=
EX^2-(EX)^2
=
\frac{(b-a)^2}{12}
\]

\end{proofbox}

\begin{proofbox}{свойства 2}

Пусть
\[
Z=\frac{X-\mu}{\sigma}
\sim\mathcal{N}(0,1)
\]

По симметрии стандартного нормального распределения
\[
EZ=0
\]

Кроме того,
\[
EZ^2=1
\]

Так как
\[
X=\mu+\sigma Z
\]
то
\[
EX=\mu+\sigma EZ=\mu
\]
и
\[
\operatorname{Var}X
=
\sigma^2\operatorname{Var}Z
=
\sigma^2
\]

\end{proofbox}

\begin{proofbox}{свойства 3}

Для экспоненциального распределения
\[
EX
=
\int_0^{+\infty}
x\lambda e^{-\lambda x}\,dx
=
\frac1\lambda
\]

Также
\[
EX^2
=
\int_0^{+\infty}
x^2\lambda e^{-\lambda x}\,dx
=
\frac{2}{\lambda^2}
\]

Следовательно,
\[
\operatorname{Var}X
=
EX^2-(EX)^2
=
\frac1{\lambda^2}
\]

\end{proofbox}

\begin{proofbox}{свойства 4}

Используем равенство
\[
\Gamma(\alpha+1)=\alpha\Gamma(\alpha)
\]

Тогда
\[
\begin{aligned}
EX
&=
\frac{\lambda^\alpha}{\Gamma(\alpha)}
\int_0^{+\infty}
x^\alpha e^{-\lambda x}\,dx\\
&=
\frac{\Gamma(\alpha+1)}
{\lambda\Gamma(\alpha)}
=
\frac{\alpha}{\lambda}
\end{aligned}
\]

Аналогично
\[
EX^2
=
\frac{\Gamma(\alpha+2)}
{\lambda^2\Gamma(\alpha)}
=
\frac{\alpha(\alpha+1)}{\lambda^2}
\]

Поэтому
\[
\operatorname{Var}X
=
EX^2-(EX)^2
=
\frac{\alpha}{\lambda^2}
\]

\end{proofbox}

\begin{proofbox}{свойства 5}

Для стандартного распределения Коши
\[
p_X(x)
=
\frac{1}{\pi(1+x^2)}
\]

Для существования математического ожидания требуется абсолютная сходимость
\[
\int_{\mathbb{R}}
|x|p_X(x)\,dx
<
+\infty
\]

Однако
\[
\begin{aligned}
\int_{\mathbb{R}}
\frac{|x|}{\pi(1+x^2)}\,dx
&=
\frac{2}{\pi}
\int_0^{+\infty}
\frac{x}{1+x^2}\,dx\\
&=
+\infty
\end{aligned}
\]

Следовательно, математическое ожидание стандартного распределения Коши не существует

\end{proofbox}

\begin{definition}

\textbf{Модой} случайной величины $X$ с плотностью $p_X$ называется точка максимума плотности:
\[
\operatorname{mod}X
=
\operatorname*{arg\,max}_{x\in\mathbb{R}}p_X(x)
\]

\end{definition}

\begin{definition}

\textbf{$\alpha$-квантилью} случайной величины $X$ называется число $q_\alpha$, для которого
\[
P(X\leq q_\alpha)\geq\alpha,
\qquad
P(X\geq q_\alpha)\geq1-\alpha
\]

Если функция распределения непрерывна и строго возрастает, то
\[
F_X(q_\alpha)=\alpha
\]

\end{definition}

\begin{definition}

\textbf{Медианой} случайной величины $X$ называется число, минимизирующее среднее абсолютное отклонение:
\[
\operatorname{med}X
=
\operatorname*{arg\,min}_{c\in\mathbb{R}}E|X-c|
\]

\end{definition}

\begin{definition}

Если $\sigma^2=\operatorname{Var}X>0$, то \textbf{коэффициентом асимметрии} называется
\[
\gamma_1
=
\frac{E(X-EX)^3}{\sigma^3}
\]

При $\gamma_1=0$ распределение симметрично относительно среднего. Знак коэффициента характеризует направление асимметрии распределения

\end{definition}

\begin{definition}

\textbf{Коэффициентом эксцесса} называется
\[
\gamma_2
=
\frac{E(X-EX)^4}{\sigma^4}-3
\]

Для нормального распределения
\[
\gamma_2=0
\]

Положительный эксцесс соответствует более острой вершине распределения, отрицательный --- более плоской

\end{definition}

\begin{definition}

\textbf{Ковариацией} случайных величин $X$ и $Y$ называется
\[
\operatorname{Cov}(X,Y)
=
E\bigl[(X-EX)(Y-EY)\bigr]
\]

Для совместной плотности $p_{X,Y}$
\[
\operatorname{Cov}(X,Y)
=
\iint_{\mathbb{R}^2}
(x-EX)(y-EY)p_{X,Y}(x,y)\,dx\,dy
\]

Также
\[
\operatorname{Cov}(X,Y)
=
E(XY)-EX\,EY
\]

\end{definition}

\begin{props}{ковариации}

1) Симметричность:
\[
\operatorname{Cov}(X,Y)
=
\operatorname{Cov}(Y,X)
\]

\medskip

2)
\[
\operatorname{Cov}(X,X)
=
\operatorname{Var}X
\]

\medskip

3) Линейность по каждому аргументу:
\[
\operatorname{Cov}(a_1X_1+a_2X_2,Y)
=
a_1\operatorname{Cov}(X_1,Y)
+
a_2\operatorname{Cov}(X_2,Y)
\]

\medskip

4) Если $X$ и $Y$ независимы, то
\[
\operatorname{Cov}(X,Y)=0
\]

Обратное утверждение в общем случае неверно

\end{props}

\begin{proofbox}{свойства 1}

По определению
\[
\begin{aligned}
\operatorname{Cov}(X,Y)
&=
E[(X-EX)(Y-EY)]\\
&=
E[(Y-EY)(X-EX)]\\
&=
\operatorname{Cov}(Y,X)
\end{aligned}
\]

\end{proofbox}

\begin{proofbox}{свойства 2}

Имеем
\[
\operatorname{Cov}(X,X)
=
E(X-EX)^2
=
\operatorname{Var}X
\]

\end{proofbox}

\begin{proofbox}{свойства 3}

Используя линейность математического ожидания, получаем
\[
\begin{aligned}
&\operatorname{Cov}(a_1X_1+a_2X_2,Y)\\
&=
E\left[
\left(
a_1(X_1-EX_1)+a_2(X_2-EX_2)
\right)(Y-EY)
\right]\\
&=
a_1\operatorname{Cov}(X_1,Y)
+
a_2\operatorname{Cov}(X_2,Y)
\end{aligned}
\]

\end{proofbox}

\begin{proofbox}{свойства 4}

Если $X$ и $Y$ независимы, то
\[
E(XY)=EX\,EY
\]

Следовательно,
\[
\operatorname{Cov}(X,Y)
=
E(XY)-EX\,EY
=
0
\]

\end{proofbox}

Из неравенства Коши--Буняковского
\[
|\operatorname{Cov}(X,Y)|
\leq
\sqrt{\operatorname{Var}X\,\operatorname{Var}Y}
\]

\begin{definition}

\textbf{Коэффициентом корреляции} случайных величин с ненулевыми дисперсиями называется
\[
\rho(X,Y)
=
\frac{\operatorname{Cov}(X,Y)}
{\sqrt{\operatorname{Var}X\,\operatorname{Var}Y}}
\]

\end{definition}

\begin{props}{коэффициента корреляции}

1)
\[
|\rho(X,Y)|\leq1
\]

\medskip

2) Равенство
\[
|\rho(X,Y)|=1
\]
достигается тогда и только тогда, когда
\[
Y=aX+b
\]
почти наверное для некоторых $a\neq0$ и $b\in\mathbb{R}$

\end{props}

\begin{proofbox}{свойства 1}

Применим неравенство Коши--Буняковского к случайным величинам
\[
X-EX,
\qquad
Y-EY
\]

Получаем
\[
\begin{aligned}
|\operatorname{Cov}(X,Y)|
&=
|E[(X-EX)(Y-EY)]|\\
&\leq
\sqrt{E(X-EX)^2}
\sqrt{E(Y-EY)^2}\\
&=
\sqrt{\operatorname{Var}X\,\operatorname{Var}Y}
\end{aligned}
\]

Деление на
\[
\sqrt{\operatorname{Var}X\,\operatorname{Var}Y}>0
\]
даёт
\[
|\rho(X,Y)|\leq1
\]

\end{proofbox}

\begin{proofbox}{свойства 2}

Равенство в неравенстве Коши--Буняковского достигается тогда и только тогда, когда случайные величины $X-EX$ и $Y-EY$ линейно зависимы почти наверное

Следовательно, существует $a\neq0$ такое, что
\[
Y-EY=a(X-EX)
\quad\text{почти наверное}
\]

Отсюда
\[
Y=aX+b
\quad\text{почти наверное}
\]
где
\[
b=EY-aEX
\]

Обратно, если
\[
Y=aX+b
\quad\text{почти наверное}
\]
то
\[
\operatorname{Cov}(X,Y)
=
a\operatorname{Var}X
\]
и
\[
\operatorname{Var}Y
=
a^2\operatorname{Var}X
\]

Поэтому
\[
|\rho(X,Y)|=1
\]

\end{proofbox}

\subsection{Числовые характеристики случайного вектора}

\begin{definition}

Математическое ожидание случайного вектора
\[
X=(X_1,\ldots,X_n)^T
\]
определяется покомпонентно:
\[
EX
=
(EX_1,\ldots,EX_n)^T
\]

\end{definition}

\begin{props}{математического ожидания случайного вектора}

Для матрицы $A$ и вектора $b$
\[
E(AX+b)=AEX+b
\]

\end{props}

\begin{proofbox}{}

Каждая компонента вектора $AX+b$ является линейной комбинацией компонент $X$. Поэтому из линейности математического ожидания
\[
\begin{aligned}
E(AX+b)_i
&=
E\left(
\sum_j A_{ij}X_j+b_i
\right)\\
&=
\sum_jA_{ij}EX_j+b_i
\end{aligned}
\]

Следовательно,
\[
E(AX+b)=AEX+b
\]

\end{proofbox}

\begin{definition}

\textbf{Ковариационной матрицей} случайного вектора $X$ называется
\[
\operatorname{Var}X
=
\operatorname{Cov}(X)
=
E\bigl[(X-EX)(X-EX)^T\bigr]
\]

Элементы матрицы ковариаций равны
\[
\operatorname{Cov}(X)_{ij}
=
\operatorname{Cov}(X_i,X_j)
\]

\end{definition}

\begin{props}{ковариационной матрицы}

1) Ковариационная матрица симметрична:
\[
\operatorname{Cov}(X)
=
\operatorname{Cov}(X)^T
\]

\medskip

2) Ковариационная матрица неотрицательно определена:
\[
t^T\operatorname{Cov}(X)t\geq0
\qquad
\text{для всех }t\in\mathbb{R}^n
\]

\medskip

3) Для линейного преобразования
\[
\operatorname{Cov}(AX+b)
=
A\operatorname{Cov}(X)A^T
\]

\medskip

4)
\[
\operatorname{Var}(X\pm Y)
=
\operatorname{Var}X+\operatorname{Var}Y
\pm2\operatorname{Cov}(X,Y)
\]

\end{props}

\begin{proofbox}{свойства 1}

Так как
\[
\operatorname{Cov}(X_i,X_j)
=
\operatorname{Cov}(X_j,X_i)
\]
то
\[
\operatorname{Cov}(X)_{ij}
=
\operatorname{Cov}(X)_{ji}
\]

Следовательно,
\[
\operatorname{Cov}(X)
=
\operatorname{Cov}(X)^T
\]

\end{proofbox}

\begin{proofbox}{свойства 2}

Для любого $t\in\mathbb{R}^n$
\[
\begin{aligned}
t^T\operatorname{Cov}(X)t
&=
t^TE[(X-EX)(X-EX)^T]t\\
&=
E\left[t^T(X-EX)(X-EX)^Tt\right]\\
&=
E\left[(t^T(X-EX))^2\right]\\
&\geq0
\end{aligned}
\]

\end{proofbox}

\begin{proofbox}{свойства 3}

Имеем
\[
E(AX+b)=AEX+b
\]

Следовательно,
\[
\begin{aligned}
\operatorname{Cov}(AX+b)
&=
E\left[
(AX-AEX)(AX-AEX)^T
\right]\\
&=
E\left[
A(X-EX)(X-EX)^TA^T
\right]\\
&=
A\operatorname{Cov}(X)A^T
\end{aligned}
\]

\end{proofbox}

\begin{proofbox}{свойства 4}

По формуле для дисперсии линейной комбинации
\[
\begin{aligned}
\operatorname{Var}(X\pm Y)
&=
E\left[
(X-EX)\pm(Y-EY)
\right]^2\\
&=
\operatorname{Var}X+\operatorname{Var}Y
\pm2\operatorname{Cov}(X,Y)
\end{aligned}
\]

\end{proofbox}

\begin{example}

Стандартный гауссовский случайный вектор
\[
X=(X_1,\ldots,X_n)^T
\]
имеет плотность
\[
p_X(x)
=
\prod_{i=1}^n
\frac{1}{\sqrt{2\pi}}e^{-x_i^2/2}
=
\frac{1}{(2\pi)^{n/2}}e^{-x^Tx/2}
\]

Поэтому его компоненты независимы и имеют распределение $\mathcal{N}(0,1)$, а
\[
EX=0,
\qquad
\operatorname{Cov}(X)=I_n
\]

\medskip

Пусть
\[
Y=AX+b
\]

Тогда
\[
Y\sim\mathcal{N}(b,AA^T)
\]
и
\[
EY=AEX+b=b
\]
\[
\operatorname{Cov}(Y)
=
A\operatorname{Cov}(X)A^T
=
AA^T
\]

\medskip

Пусть
\[
X\sim\mathcal{N}(\mu,\Sigma)
\]
где
\[
\Sigma=\Sigma^T,
\qquad
\Sigma\geq0
\]

Если $Z\sim\mathcal{N}(0,I_n)$, то можно представить
\[
X=\sqrt{\Sigma}\,Z+\mu
\]

Тогда
\[
EX=\mu
\]
и
\[
\operatorname{Cov}(X)
=
\sqrt{\Sigma}\,I_n\,(\sqrt{\Sigma})^T
=
\Sigma
\]

\medskip

Если ковариационная матрица нормального вектора диагональна,
\[
X\sim
\mathcal{N}
\left(
\mu,
\operatorname{diag}(\lambda_1,\ldots,\lambda_n)
\right)
\]
и все $\lambda_i>0$, то совместная плотность равна
\[
p_X(x)
=
\prod_{i=1}^n
\frac{1}{\sqrt{2\pi\lambda_i}}
\exp\left(
-\frac{(x_i-\mu_i)^2}{2\lambda_i}
\right)
\]

Следовательно, компоненты $X_1,\ldots,X_n$ независимы

\end{example}

\newpage

\section{Условные распределения и числовые характеристики}

\subsection{Условное распределение}

\begin{definition}

Для дискретных случайных величин $X$ и $Y$ условное распределение $X$ при условии $Y=y_j$ задаётся равенством
\[
P(X=x_i\mid Y=y_j)
=
\frac{P(X=x_i,Y=y_j)}
{P(Y=y_j)}
\]
при
\[
P(Y=y_j)>0
\]

\end{definition}

\begin{definition}

Для абсолютно непрерывного случайного вектора $(X,Y)$ \textbf{условной плотностью} $X$ при условии $Y=y$ называется
\[
p_{X\mid Y}(x\mid y)
=
\frac{p_{X,Y}(x,y)}{p_Y(y)}
\]
при
\[
p_Y(y)>0
\]

\end{definition}

Рассмотрим условие
\[
y\leq Y\leq y+\Delta y
\]

Тогда
\[
P(X\leq x\mid y\leq Y\leq y+\Delta y)
=
\frac{
P(X\leq x,\ y\leq Y\leq y+\Delta y)
}{
P(y\leq Y\leq y+\Delta y)
}
\]

Для абсолютно непрерывного распределения
\[
P(X\leq x\mid y\leq Y\leq y+\Delta y)
=
\frac{
\displaystyle
\int_{-\infty}^x
\int_y^{y+\Delta y}
p(u,v)\,dv\,du
}{
\displaystyle
\int_y^{y+\Delta y}p_Y(v)\,dv
}
\]

При $\Delta y\to0$ получаем
\[
F_{X\mid Y}(x\mid y)
=
\int_{-\infty}^x
\frac{p(u,y)}{p_Y(y)}\,du
\]

Следовательно,
\[
p_{X\mid Y}(x\mid y)
=
\frac{p(x,y)}{p_Y(y)}
\]

Также
\[
p_X(x)
=
\int_{-\infty}^{+\infty}
p_{X\mid Y}(x\mid y)p_Y(y)\,dy
\]

\begin{theorem}{Байеса для непрерывного случая}

Если соответствующие плотности положительны, то
\[
p_{X\mid Y}(x\mid y)
=
\frac{
p_{Y\mid X}(y\mid x)p_X(x)
}{
p_Y(y)
}
\]

\end{theorem}

\begin{proofbox}{}

По определению условной плотности
\[
p_{X\mid Y}(x\mid y)
=
\frac{p_{X,Y}(x,y)}{p_Y(y)}
\]

С другой стороны,
\[
p_{X,Y}(x,y)
=
p_{Y\mid X}(y\mid x)p_X(x)
\]

Следовательно,
\[
p_{X\mid Y}(x\mid y)
=
\frac{
p_{Y\mid X}(y\mid x)p_X(x)
}{
p_Y(y)
}
\]

\end{proofbox}

\begin{example}

Пусть
\[
X=(X_1,X_2,X_3)^T
\sim
\mathcal{N}(0,I_3)
\]
и
\[
Y
=
\frac{X_1+X_2X_3}{\sqrt{1+X_3^2}}
\]

Найдём распределение $Y$

Зафиксируем $X_3=x_3$. Тогда
\[
Y\mid X_3=x_3
=
\frac{X_1+x_3X_2}{\sqrt{1+x_3^2}}
\]

Так как $X_1$ и $X_2$ независимы и имеют распределение $\mathcal{N}(0,1)$, то
\[
X_1+x_3X_2
\sim
\mathcal{N}(0,1+x_3^2)
\]

Следовательно,
\[
Y\mid X_3=x_3
\sim
\mathcal{N}(0,1)
\]

Условное распределение не зависит от $x_3$, поэтому
\[
p_Y(y)
=
\int_{-\infty}^{+\infty}
p_{Y\mid X_3}(y\mid x_3)p_{X_3}(x_3)\,dx_3
\]

Отсюда
\[
\begin{aligned}
p_Y(y)
&=
\int_{-\infty}^{+\infty}
\frac{1}{\sqrt{2\pi}}e^{-y^2/2}
\frac{1}{\sqrt{2\pi}}e^{-x_3^2/2}\,dx_3\\
&=
\frac{1}{\sqrt{2\pi}}e^{-y^2/2}
\end{aligned}
\]

Таким образом,
\[
Y\sim\mathcal{N}(0,1)
\]

\end{example}

\subsection{Условное математическое ожидание}

\begin{definition}

Для абсолютно непрерывных случайных величин $X$ и $Y$ \textbf{условным математическим ожиданием} $X$ при условии $Y=y$ называется
\[
E(X\mid Y=y)
=
\int_{-\infty}^{+\infty}
x\,p_{X\mid Y}(x\mid y)\,dx
\]

Величина
\[
E(X\mid Y)
\]
является случайной величиной, являющейся функцией от $Y$

\end{definition}

\begin{props}{условного математического ожидания}

1) Линейность:
\[
E(a_1X_1+a_2X_2\mid Y)
=
a_1E(X_1\mid Y)+a_2E(X_2\mid Y)
\]

\medskip

2) Если $f(Y)$ интегрируема, то $E(f(Y)\mid Y)=f(Y)$

\medskip

3) Если $X$ и $Y$ независимы, то $E(X\mid Y)=EX$

\medskip

4) Формула полного математического ожидания:
\[
EX
=
E\bigl[E(X\mid Y)\bigr]
\]

\medskip

5) Условное математическое ожидание является наилучшим среднеквадратичным приближением $X$ функцией от $Y$:
\[
E(X\mid Y)
=
\operatorname*{arg\,min}_{g}
E(X-g(Y))^2
\]

\end{props}

\begin{proofbox}{свойства 1}

Из линейности интеграла
\[
\begin{aligned}
&E(a_1X_1+a_2X_2\mid Y=y)\\
&=
a_1E(X_1\mid Y=y)
+
a_2E(X_2\mid Y=y)
\end{aligned}
\]

Следовательно,
\[
E(a_1X_1+a_2X_2\mid Y)
=
a_1E(X_1\mid Y)+a_2E(X_2\mid Y)
\]

\end{proofbox}

\begin{proofbox}{Доказательство свойства 2}

При фиксированном значении $Y=y$ величина $f(Y)$ равна константе $f(y)$

Поэтому
\[
E(f(Y)\mid Y=y)
=
f(y)
\]

Следовательно,
\[
E(f(Y)\mid Y)=f(Y)
\]

\end{proofbox}

\begin{proofbox}{свойства 3}

Если $X$ и $Y$ независимы, то
\[
p_{X\mid Y}(x\mid y)=p_X(x)
\]

Поэтому
\[
\begin{aligned}
E(X\mid Y=y)
&=
\int_{\mathbb{R}}
x p_{X\mid Y}(x\mid y)\,dx\\
&=
\int_{\mathbb{R}}
x p_X(x)\,dx\\
&=
EX
\end{aligned}
\]

Следовательно,
\[
E(X\mid Y)=EX
\]

\end{proofbox}

\begin{proofbox}{свойства 4}

В абсолютно непрерывном случае
\[
\begin{aligned}
EX
&=
\int_{-\infty}^{+\infty}
x p_X(x)\,dx\\
&=
\int_{-\infty}^{+\infty}
x
\left(
\int_{-\infty}^{+\infty}
p_{X\mid Y}(x\mid y)p_Y(y)\,dy
\right)dx\\
&=
\int_{-\infty}^{+\infty}
\left(
\int_{-\infty}^{+\infty}
x p_{X\mid Y}(x\mid y)\,dx
\right)p_Y(y)\,dy\\
&=
\int_{-\infty}^{+\infty}
E(X\mid Y=y)p_Y(y)\,dy\\
&=
E[E(X\mid Y)]
\end{aligned}
\]

\end{proofbox}

\begin{proofbox}{свойства 5}

Пусть $g(Y)$ --- произвольная функция от $Y$. Представим
\[
X-g(Y)
=
X-E(X\mid Y)+E(X\mid Y)-g(Y)
\]

Тогда
\[
\begin{aligned}
E(X-g(Y))^2
&=
E(X-E(X\mid Y))^2\\
&\quad+
E(E(X\mid Y)-g(Y))^2\\
&\quad+
2E\left[
(X-E(X\mid Y))
(E(X\mid Y)-g(Y))
\right]
\end{aligned}
\]

Для последнего слагаемого по формуле полного математического ожидания
\[
\begin{aligned}
&E\left[
(X-E(X\mid Y))
(E(X\mid Y)-g(Y))
\right]\\
&=
E\left[
E\left(
(X-E(X\mid Y))
(E(X\mid Y)-g(Y))
\mid Y
\right)
\right]\\
&=
E\left[
(E(X\mid Y)-g(Y))
E(X-E(X\mid Y)\mid Y)
\right]\\
&=
0
\end{aligned}
\]

Следовательно,
\[
E(X-g(Y))^2
=
E(X-E(X\mid Y))^2
+
E(E(X\mid Y)-g(Y))^2
\]

Минимум достигается при
\[
g(Y)=E(X\mid Y)
\]

\end{proofbox}

\subsection{Геометрическая интерпретация условного математического ожидания}

Рассмотрим пространство случайных величин с конечным вторым моментом и введём скалярное произведение
\[
\langle X,Y\rangle
=
E(XY)
\]

Оно линейно по каждому аргументу и удовлетворяет
\[
\langle X,X\rangle
=
EX^2\geq0
\]

Причём
\[
EX^2=0
\quad\Longrightarrow\quad
X=0
\quad\text{почти наверное}
\]

Соответствующая норма имеет вид
\[
\|X\|
=
\sqrt{EX^2}
\]
а расстояние между случайными величинами равно
\[
\rho(X,Y)
=
\|X-Y\|
\]

Пусть
\[
\mathcal{H}_Y
=
\left\{
g(Y):Eg^2(Y)<+\infty
\right\}
\]

Тогда
\[
E(X\mid Y)
\]
является ортогональной проекцией $X$ на пространство $\mathcal{H}_Y$

Действительно, для любой $g(Y)\in\mathcal{H}_Y$
\[
E\left[
(X-E(X\mid Y))g(Y)
\right]
=
0
\]

\subsection{Условная дисперсия}

\begin{definition}

\textbf{Условной дисперсией} случайной величины $X$ при условии $Y$ называется
\[
\operatorname{Var}(X\mid Y)
=
E\left[
(X-E(X\mid Y))^2
\mid Y
\right]
\]

Эквивалентно,
\[
\operatorname{Var}(X\mid Y)
=
E(X^2\mid Y)-E^2(X\mid Y)
\]

\end{definition}

\begin{theorem}{о полной дисперсии}

\[
\operatorname{Var}X
=
\operatorname{Var}(E(X\mid Y))
+
E[\operatorname{Var}(X\mid Y)]
\]

\end{theorem}

\begin{proofbox}{}

\[
\begin{aligned}
\operatorname{Var}X
&=
EX^2-(EX)^2\\
&=
E\bigl(E(X^2\mid Y)\bigr)
-
\bigl(E(E(X\mid Y))\bigr)^2\\
&=
E\bigl(\operatorname{Var}(X\mid Y)+E^2(X\mid Y)\bigr)
-
\bigl(E(E(X\mid Y))\bigr)^2\\
&=
E\operatorname{Var}(X\mid Y)
+
\operatorname{Var}(E(X\mid Y))
\end{aligned}
\]

\end{proofbox}

\begin{example}

Пусть
\[
\begin{pmatrix}
X_1\\
X_2
\end{pmatrix}
\sim
\mathcal{N}
\left(
\begin{pmatrix}
\mu_1\\
\mu_2
\end{pmatrix},
\begin{pmatrix}
\Sigma_{11} & \Sigma_{12}\\
\Sigma_{21} & \Sigma_{22}
\end{pmatrix}
\right)
\]

Найдём условное распределение $X_1$ при фиксированном $X_2$

Положим
\[
U=X_1-AX_2
\]

Выберем $A$ так, чтобы $U$ и $X_2$ были некоррелированы:
\[
\operatorname{Cov}(U,X_2)=0
\]

Имеем
\[
\begin{aligned}
\operatorname{Cov}(X_1-AX_2,X_2)
&=
\Sigma_{12}-A\Sigma_{22}
\end{aligned} = 0
\]

Не умаляя общности, предположим, что компоненты $X_2$ невырожденные, то есть
\[
\Sigma_{22}>0
\]

Докажем, что из нулевой компоненты $X_2$ следует
\[
\Sigma_{22}>0
\]

Предположим противное, то есть
\[
\Sigma_{22}\geq 0
\]
и существует $t\neq 0$ такое, что
\[
t^T\Sigma_{22}t=0
\]

Тогда
\[
\underbrace{\operatorname{Var}(t^TX_2)}_{=0}=0
\]
следовательно,
\[
t^TX_2=\operatorname{const}
\]
то есть компоненты $X_2$ линейно зависимы, что приводит к противоречию

Рассмотрим
\[
U=X_1-AX_2
\]

Подберём $A$ так, чтобы $U$ и $X_2$ были некоррелированы:
\[
\begin{aligned}
\operatorname{Cov}(X_1-AX_2,X_2)
&=
\operatorname{Cov}(X_1,X_2)
-
A\operatorname{Cov}(X_2,X_2)\\
&=
\Sigma_{12}-A\Sigma_{22}
=
0
\end{aligned}
\]

Отсюда
\[
A=\Sigma_{12}\Sigma_{22}^{-1}
\]

Имеем
\[
X_1=U+AX_2
\]

Тогда распределение $X_1$ при условии $X_2$ имеет вид
\[
X_1\mid X_2
\sim
\mathcal{N}
\left(
E(X_1\mid X_2),
\operatorname{Var}(X_1\mid X_2)
\right)
\]

Найдём условное математическое ожидание:
\[
\begin{aligned}
E(X_1\mid X_2)
&=
E(U+AX_2\mid X_2)\\
&=
EU+AX_2\\
&=
\mu_1-A\mu_2+AX_2\\
&=
\mu_1+A(X_2-\mu_2)\\
&=
\mu_1+
\Sigma_{12}\Sigma_{22}^{-1}(X_2-\mu_2)
\end{aligned}
\]

Найдём условную дисперсию:
\[
\begin{aligned}
\operatorname{Var}(X_1\mid X_2)
&=
\operatorname{Var}(U+AX_2\mid X_2)\\
&=
\operatorname{Var}U
+
A\operatorname{Var}(X_2\mid X_2)A^T\\
&=
\operatorname{Var}U\\
&=
\operatorname{Var}(X_1-AX_2)\\
&=
\operatorname{Cov}(X_1-AX_2,X_1-AX_2)\\
&=
\operatorname{Var}X_1
-
A\operatorname{Cov}(X_2,X_1)
-
\operatorname{Cov}(X_1,X_2)A^T
+
A\operatorname{Cov}(X_2,X_2)A^T\\
&=
\Sigma_{11}
-
\Sigma_{12}\Sigma_{22}^{-1}\Sigma_{21}
-
\Sigma_{12}\Sigma_{22}^{-1}\Sigma_{21}
+
\Sigma_{12}\Sigma_{22}^{-1}
\Sigma_{22}
\Sigma_{22}^{-1}\Sigma_{21}\\
&=
\Sigma_{11}
-
\Sigma_{12}\Sigma_{22}^{-1}\Sigma_{21}
\end{aligned}
\]

\begin{corollary}

$X_1,X_2$ --- 1d величины

\[
E(X_1\mid X_2)
=
\mu_1
+
\operatorname{cov}(X_1,X_2)\cdot
\sigma_2^{-2}(X_2-\mu_2)
\]

\[
\operatorname{Var}(X_1\mid X_2)
=
\sigma_1^2
-
\frac{\operatorname{cov}^2(X_1,X_2)}
{\sigma_2^2}
\]

\end{corollary}

\end{example}

\newpage

\section{Сходимости случайных величин}

\begin{definition}

\textbf{Сходимость по вероятности}

Пусть $(\Omega,\mathcal{A},P)$ --- вероятностное пространство, $X_n,X$ --- случайные величины на $(\Omega,\mathcal{A},P)$

Говорят, что $X_n$ сходится к $X$ по вероятности, если
\[
X_n\xrightarrow[n\to\infty]{P}X
\quad\Longleftrightarrow\quad
P(|X_n-X|\geq\varepsilon)\xrightarrow[n\to\infty]{}0
\\
\Longleftrightarrow P(|X_n-X|<\varepsilon)\xrightarrow[n\to\infty]{}1
\]
для любого $\varepsilon>0$

\end{definition}

\begin{definition}

\textbf{Сходимость по распределению}

Говорят, что $X_n$ сходится к $X$ по распределению, если
\[
X_n\xrightarrow[n\to\infty]{d}X
\quad\Longleftrightarrow\quad
F_n(t)\xrightarrow[n\to\infty]{}F(t)
\qquad
\forall t\in C(F)
\]
где
\[
C(F)=\{t:F\text{ непрерывна в }t\}
\]

Также используется обозначение
\[
F_n\xrightarrow{d}F
\]

\end{definition}

\begin{note}
1) Сходимость по распределению, вообще говоря, не влечёт сходимость по вероятности

\[
\xrightarrow{d} \ \not\Rightarrow \ \xrightarrow{P}
\]

\bigskip

2) Если $F\in C(\mathbb{R})$, то
   \[
   X_n\xrightarrow{d}X
   \quad\Longleftrightarrow\quad
   F_n(t)\longrightarrow F(t)
   \quad \forall t\in\mathbb{R},
   \]
   где $F_n$ и $F$ --- функции распределения $X_n$ и $X$ соответственно.

\bigskip

3)  $X_n, X$ - дискретные, $\sup pX_n = \sup pX$. Тогда
\[
X_n \xrightarrow{d}X \Longleftrightarrow P(X_n=t_m) \to P(X = t_m) \ \forall t_m 
\]

\end{note}

\begin{proofbox}{пункта 1}
Пусть
\[
P\left(X=1\right)=P\left(X=-1\right)=\frac12,
\qquad
X_n=(-1)^nX.
\]
Тогда для любого $n$
\[
X_n\overset{d}{=}X,
\]
поэтому
\[
X_n\xrightarrow{d}X.
\]
Однако
\[
|X_n-X|=
\begin{cases}
0, & n \text{ чётное},\\
2, & n \text{ нечётное}.
\end{cases}
\]
Следовательно, при $0<\varepsilon<2$
\[
P\left(|X_n-X|>\varepsilon\right)=
\begin{cases}
0, & n \text{ чётное},\\
1, & n \text{ нечётное},
\end{cases}
\]
и потому $X_n\not\xrightarrow{p}X$.
\end{proofbox}

\begin{proofbox}{пункта 2}
По определению сходимости по распределению
\[
X_n\xrightarrow{d}X
\]
означает, что
\[
F_n(t)\longrightarrow F(t)
\]
во всех точках непрерывности функции распределения $F$.

Так как $F\in C(\mathbb{R})$, каждая точка $t\in\mathbb{R}$ является точкой непрерывности $F$. Поэтому
\[
X_n\xrightarrow{d}X
\quad\Longleftrightarrow\quad
F_n(t)\longrightarrow F(t)
\quad\forall t\in\mathbb{R}.
\]
\end{proofbox}

\begin{proofbox}{пункта 3}
Обозначим общее дискретное множество значений
\[
\operatorname{supp}X_n=\operatorname{supp}X=\{t_1,t_2,\ldots\}.
\]

\(\Rightarrow\)

Пусть
\[
X_n\xrightarrow{d}X
\]
Тогда для любой точки непрерывности \(x\) функции распределения \(F_X\)
\[
F_{X_n}(x)\to F_X(x)
\]

Зафиксируем \(t_m\). Поскольку множество \(\{t_k\}\) дискретно, можно выбрать \(a<t_m<b\) так, чтобы
\[
(a,b)\cap\operatorname{supp}X=\{t_m\}
\]
причём \(a,b\notin\operatorname{supp}X\). Следовательно, \(a,b\) — точки непрерывности \(F_X\).

Так как носители \(X_n\) и \(X\) совпадают,
\[
P(X_n=t_m)=F_{X_n}(b)-F_{X_n}(a).
\]
По сходимости по распределению
\[
F_{X_n}(b)\to F_X(b),\qquad
F_{X_n}(a)\to F_X(a).
\]
Значит,
\[
P(X_n=t_m)
\to F_X(b)-F_X(a)
=P(X=t_m).
\]

Таким образом,
\[
P(X_n=t_m)\to P(X=t_m)
\qquad\forall t_m.
\]

\(\Leftarrow\)

Пусть теперь
\[
P(X_n=t_m)\to P(X=t_m)
\qquad\forall m.
\]

Нужно доказать
\[
X_n\xrightarrow{d}X
\]

Возьмём произвольную точку непрерывности \(x\) функции \(F_X\). Имеем
\[
F_{X_n}(x)
=
\sum_{t_m\le x}P(X_n=t_m),
\qquad
F_X(x)
=
\sum_{t_m\le x}P(X=t_m).
\]

Если множество значений конечно, результат сразу следует из почленной сходимости конечной суммы.

В общем счётном случае зафиксируем \(\varepsilon>0\). Поскольку
\[
\sum_{m=1}^{\infty}P(X=t_m)=1,
\]
можно выбрать конечное множество индексов \(I\), для которого
\[
\sum_{m\notin I}P(X=t_m)<\varepsilon.
\]
Из
\[
P(X_n=t_m)\to P(X=t_m),\qquad m\in I,
\]
следует
\[
\sum_{m\in I}P(X_n=t_m)
\to
\sum_{m\in I}P(X=t_m).
\]
Поэтому при достаточно больших \(n\)
\[
\sum_{m\notin I}P(X_n=t_m)<2\varepsilon.
\]

Теперь
\[
\begin{aligned}
|F_{X_n}(x)-F_X(x)|
&\le
\sum_{\substack{m\in I\\t_m\le x}}
|P(X_n=t_m)-P(X=t_m)|
\\
&\quad+
\sum_{m\notin I}P(X_n=t_m)
+
\sum_{m\notin I}P(X=t_m).
\end{aligned}
\]
Первое слагаемое стремится к нулю, а два остальных при больших \(n\) меньше \(2\varepsilon\) и \(\varepsilon\). Поскольку \(\varepsilon>0\) произвольно,
\[
F_{X_n}(x)\to F_X(x).
\]
Следовательно,
\[
X_n\xrightarrow{d}X.
\]

Итак,
\[
X_n\xrightarrow{d}X
\Longleftrightarrow
P(X_n=t_m)\to P(X=t_m)
\quad\forall t_m.
\]

Здесь существенно условие об общем дискретном носителе: оно позволяет в первой части изолировать каждую точку \(t_m\) одинаково для \(X_n\) и \(X\)
\end{proofbox}

\begin{props}{сходимостей случайных величин}

1) Если $X_n\xrightarrow{d}c$, то $X_n\xrightarrow{p}c$

\bigskip

2) \textbf{Теорема Слуцкого}

Если $X_n\xrightarrow{d}X, \ \ Y_n\xrightarrow{p}c$,
   где $c\in\mathbb{R}$ --- константа, то $X_n + Y_n\xrightarrow{d}X + c$

\bigskip

3) Если $X_n\xrightarrow{d}X, \ \ Y_n\xrightarrow{p}c$,
   где $c\in\mathbb{R}$ --- константа, то $X_nY_n\xrightarrow{d}cX$

\bigskip

4) Если $X_n\xrightarrow{p}X$, то $X_n\xrightarrow{d}X$

\bigskip
   
5) Пусть
   \[
   X_{1n}\xrightarrow{p}X_1,\quad
   X_{2n}\xrightarrow{p}X_2,\quad \ldots,\quad
   X_{mn}\xrightarrow{p}X_m,
   \]
   и $g\in C(\mathbb{R}^m)$. Тогда
   \[
   g(X_{1n},\ldots,X_{mn})
   \xrightarrow{p}
   g(X_1,\ldots,X_m).
   \]

\bigskip

6) Если $X_n\xrightarrow{d}X$ и $g$ непрерывна, то $g(X_n)\xrightarrow{d}g(X)$

\bigskip
   
7) Пусть
   \[
   X_n\xrightarrow{d}X,
   \qquad
   t_n\longrightarrow t\in[-\infty,+\infty].
   \]
   Тогда:
а) если $t\in\mathbb{R}$ и $t\in C(F_X)$, то
\[
F_{X_n}(t_n)\longrightarrow F_X(t);
\]
б) если $t=+\infty$, то
\[
F_{X_n}(t_n)\longrightarrow1;
\]
в) если $t=-\infty$, то
\[
F_{X_n}(t_n)\longrightarrow0.
\]
\end{props}

\begin{proofbox}{свойства 1}
\[
P(|X_n-c|<\varepsilon)
=
P\left(c-\frac{\varepsilon}{2}<X_n\leq c+\frac{\varepsilon}{2}\right)
=
F_{X_n}\left(c+\frac{\varepsilon}{2}\right)
-
F_{X_n}\left(c-\frac{\varepsilon}{2}\right).
\]
\[
\longrightarrow
F_{I(c)}\left(c+\frac{\varepsilon}{2}\right)
-
F_{I(c)}\left(c-\frac{\varepsilon}{2}\right)
=
1-0=1.
\]
\end{proofbox}

\begin{proofbox}{свойства 2}
Не умаляя общности, можно сказать, что $c=0$, так как
\[
X_n\xrightarrow{d}X
\Rightarrow
X_n-c\xrightarrow{d}X-c.
\]
\[
P(X_n+Y_n\leq t)
\overset{?}{\longrightarrow}
P(X\leq t),
\qquad
\forall t\in C(F_X).
\]
\[
P(X_n+Y_n\leq t)
=
P(X_n+Y_n\leq t,\ |Y_n|<\varepsilon)
+
P(X_n+Y_n\leq t,\ |Y_n|\geq\varepsilon).
\]
\[
P(X_n+Y_n\leq t,\ |Y_n|\geq\varepsilon)
\leq
P(|Y_n|\geq\varepsilon)
\xrightarrow[n\to\infty]{}0.
\]
\[
\begin{cases}
X_n+Y_n\leq t,\\
-\varepsilon<Y_n<\varepsilon
\end{cases}
\Rightarrow
X_n\leq t+\varepsilon.
\]
\[
P(X_n+Y_n\leq t,\ |Y_n|<\varepsilon)
\leq
P(X_n\leq t+\varepsilon).
\]
\[
P(X_n\leq t+\varepsilon)
+
P(|Y_n|\geq\varepsilon)
\Rightarrow
\overline{\lim}
P(X_n+Y_n\leq t)
\leq
F_X(t+\varepsilon).
\]
\[
P(X_n\leq t-\varepsilon)
=
P(X_n\leq t-\varepsilon,\ |Y_n|<\varepsilon)
+
P(X_n\leq t-\varepsilon,\ |Y_n|\geq\varepsilon).
\]
\[
\begin{cases}
X_n\leq t-\varepsilon,\\
-\varepsilon<Y_n<\varepsilon
\end{cases}
\Rightarrow
X_n+Y_n\leq t.
\]
\[
P(X_n\leq t-\varepsilon)
\leq
P(X_n+Y_n\leq t)
+
P(|Y_n|\geq\varepsilon).
\]
\[
\Rightarrow
P(X_n+Y_n\leq t)
\geq
P(X_n\leq t-\varepsilon)
-
P(|Y_n|\geq\varepsilon).
\]
\[
\underline{\lim}
P(X_n+Y_n\leq t)
\geq
F_X(t-\varepsilon).
\]
\end{proofbox}

\begin{proofbox}{свойства 3}
\[
X_n(Y_n-c+c)=X_n(Y_n-c)+cX_n\xrightarrow{d}cX.
\]

\[
Y_n\to0 \Rightarrow X_nY_n\xrightarrow{d}0.
\]

\[
P(X_nY_n\leq t)\longrightarrow
\begin{cases}
0, & t<0,\\
1, & t>0.
\end{cases}
\]

\[
P(-t\leq X_nY_n\leq t)\longrightarrow1
\qquad \forall t>0.
\]

\[
P(-t\leq X_nY_n\leq t)
=
P(-t\leq X_nY_n\leq t,\ |X_n|\geq M)
+
P(-t\leq X_nY_n\leq t,\ |X_n|<M).
\]

\[
P(-t\leq X_nY_n\leq t)
=
P(-t\leq X_nY_n\leq t,\ |X_n|\leq M)
+
P(-t\leq X_nY_n\leq t,\ |X_n|>M),
\]
\[
P(-t\leq X_nY_n\leq t,\ |X_n|>M)
\leq P(|X_n|>M),
\]
\[
P(|X_n|>M)\longrightarrow 1-F(M)+F(-M)<\varepsilon.
\]

\[
P(-t\leq Y_n\leq t)
=
P(-t\leq Y_n\leq t,\ |X_n|\leq M)
+
P(-t\leq Y_n\leq t,\ |X_n|>M)
\leq
P\left(-Mt\leq X_nY_n\leq Mt\right)
+
P(|X_n|>M).
\]
\end{proofbox}

\begin{proofbox}{свойства 4}
\[
X_n = (X_n - X) + X \xrightarrow{d} 0 + X = X
\]
\end{proofbox}

\begin{proofbox}{свойства 5}
\[
p(|g(X_n) - g(X)| < \varepsilon) \geq P(|X_n - X| < \varepsilon) \longrightarrow 1
\]
\end{proofbox}

\begin{proofbox}{свойства 6}
Пусть $F_{g(X)}$ --- функция распределения случайной величины $g(X)$ и
\[
t\in C(F_{g(X)}).
\]
Рассмотрим множество
\[
A_t=\{x\in\mathbb{R}:g(x)\leq t\}.
\]
Так как $g$ непрерывна, граница этого множества может содержать только точки $x$, для которых $g(x)=t$. Поэтому
\[
P(X\in\partial A_t)
\leq
P(g(X)=t).
\]
Поскольку $t$ является точкой непрерывности $F_{g(X)}$,
\[
P(g(X)=t)=0.
\]
Следовательно, $A_t$ является множеством непрерывности распределения $X$.

Из $X_n\xrightarrow{d}X$ получаем
\[
P(X_n\in A_t)
\longrightarrow
P(X\in A_t).
\]
То есть
\[
P(g(X_n)\leq t)
\longrightarrow
P(g(X)\leq t).
\]
Следовательно,
\[
g(X_n)\xrightarrow{d}g(X).
\]
\end{proofbox}

\begin{proofbox}{свойства 7}
Сначала пусть
\[
t_n\longrightarrow t\in\mathbb{R},
\qquad
t\in C(F_X).
\]
Для любого $\varepsilon>0$ при достаточно больших $n$
\[
t-\varepsilon<t_n<t+\varepsilon.
\]
Так как функции распределения неубывают,
\[
F_{X_n}(t-\varepsilon)
\leq
F_{X_n}(t_n)
\leq
F_{X_n}(t+\varepsilon).
\]
Выберем $\varepsilon>0$ так, чтобы $t-\varepsilon$ и $t+\varepsilon$ были точками непрерывности $F_X$. Из
\[
X_n\xrightarrow{d}X
\]
следует
\[
F_{X_n}(t-\varepsilon)\longrightarrow F_X(t-\varepsilon),
\qquad
F_{X_n}(t+\varepsilon)\longrightarrow F_X(t+\varepsilon).
\]
Поэтому
\[
F_X(t-\varepsilon)
\leq
\liminf_{n\to\infty}F_{X_n}(t_n)
\leq
\limsup_{n\to\infty}F_{X_n}(t_n)
\leq
F_X(t+\varepsilon).
\]
Устремляя $\varepsilon\to0$ и используя непрерывность $F_X$ в точке $t$, получаем
\[
F_{X_n}(t_n)\longrightarrow F_X(t).
\]
Пусть теперь $t_n\to+\infty$. Возьмём произвольную точку непрерывности $M$ функции $F_X$. Для достаточно больших $n$
\[
t_n\geq M,
\]
поэтому
\[
F_{X_n}(t_n)\geq F_{X_n}(M).
\]
Следовательно,
\[
\liminf_{n\to\infty}F_{X_n}(t_n)
\geq F_X(M).
\]
Устремляя $M\to+\infty$, получаем
\[
\liminf_{n\to\infty}F_{X_n}(t_n)\geq1.
\]
Так как $F_{X_n}(t_n)\leq1$, имеем
\[
F_{X_n}(t_n)\longrightarrow1.
\]
Наконец, пусть $t_n\to-\infty$. Возьмём произвольную точку непрерывности $M$ функции $F_X$. Для достаточно больших $n$
\[
t_n\leq M,
\]
следовательно,
\[
F_{X_n}(t_n)\leq F_{X_n}(M).
\]
Отсюда
\[
\limsup_{n\to\infty}F_{X_n}(t_n)
\leq F_X(M).
\]
Устремляя $M\to-\infty$, получаем
\[
\limsup_{n\to\infty}F_{X_n}(t_n)\leq0.
\]
Значит,
\[
F_{X_n}(t_n)\longrightarrow0.
\]
\end{proofbox}

\newpage

\section{Характеристические функции случайных величин}

\begin{definition}

Пусть $X$ --- случайная величина. \textbf{Характеристической функцией} случайной величины $X$ называется
\[
f_X(t)
=
Ee^{itX},
\qquad
t\in\mathbb{R}
\]

Так как
\[
e^{itX}
=
\cos(tX)+i\sin(tX)
\]
то
\[
f_X(t)
=
E\cos(tX)+iE\sin(tX)
\]

Если $X$ --- дискретная случайная величина, то
\[
f_X(t)
=
\sum_k p_k e^{itx_k}
\]

Если $X$ имеет плотность $p_X$, то
\[
f_X(t)
=
\int_{\mathbb{R}}
p_X(x)e^{itx}\,dx
\]

\end{definition}

\begin{props}{характеристической функции}

1) $f_X(0)=1 $ и $|f_X(t)|\leq1$

\bigskip

2) Если $Y=aX+b$, то $f_Y(t) = e^{itb}f_X(at)$

В частности, если $Y=-X$, то $f_{-X}(t) =
f_X(-t) = \overline{f_X(t)}$

\bigskip

3) Если $X\overset{d}{=}-X$, то $f_X(t)=\overline{f_X(t)} \Longrightarrow f_x(t) \in \mathbb{R} \ \ \forall t\in\mathbb{R}$

\bigskip

4) Если существуют моменты $EX^k$, то $f_X^{(k)}(0) = i^kEX^k$

В частности, $f_X'(0)=iEX$ и $f_X''(0)=i^2EX^2=-EX^2$

При наличии моментов до порядка $n$
\[
f_X(t)
=
1+
\sum_{k=1}^n
\frac{i^kEX^k}{k!}t^k
+
o(t^n),
\qquad
t\to0
\]

\bigskip

5) Если $X_1,\ldots,X_n$ независимы и $S_n=X_1+\cdots+X_n$, 
то $f_{S_n}(t) = \prod_{j=1}^nf_{X_j}(t)$

\bigskip

6) Характеристическая функция равномерно непрерывна:
\[
\sup_{t\in\mathbb{R}}
|f_X(t+h)-f_X(t)|
\xrightarrow[h\to0]{}0
\]

\end{props}

\begin{proofbox}{свойства 1}

Имеем
\[
f_X(0)
=
Ee^{i0X}
=
1
\]

Кроме того,
\[
|f_X(t)|
=
|Ee^{itX}|
\leq
E|e^{itX}|
=
1
\]

\end{proofbox}

\begin{proofbox}{свойства 2}

Если
\[
Y=aX+b
\]
то
\[
\begin{aligned}
f_Y(t)
&=
Ee^{itY}\\
&=
Ee^{it(aX+b)}\\
&=
e^{itb}Ee^{i(at)X}\\
&=
e^{itb}f_X(at)
\end{aligned}
\]

При $Y=-X$
\[
f_{-X}(t)
=
f_X(-t)
\]

При этом
\[
f_X(-t)
=
Ee^{-itX}
=
E\overline{e^{itX}}
=
\overline{f_X(t)}
\]

\end{proofbox}

\begin{proofbox}{свойства 3}

По предыдущему свойству $f_{-X}(t) = \overline{f_X(t)}$. Значит

\[
f_x(t) \in \mathbb{R} \ \ \forall t\in\mathbb{R}
\]

\end{proofbox}

\begin{proofbox}{свойства 4}

Дифференцируя характеристическую функцию, получаем
\[
f_X'(t)
=
\left(Ee^{itX}\right)'
=
E\left(e^{itX}\right)'
=
E\left(iXe^{itX}\right)
\]

Поэтому
\[
f_X'(0)
=
iEX
\]

Аналогично
\[
f_X''(0)
=
i^2EX^2
\]

В общем случае
\[
f_X^{(k)}(0)
=
i^kEX^k
\]

Следовательно, по формуле Тейлора
\[
f_X(t)
=
1+
\sum_{k=1}^n
\frac{i^kEX^k}{k!}t^k
+
o(t^n)
\]

\end{proofbox}

\begin{proofbox}{свойства 5}

Так как $X_1,\ldots,X_n$ независимы, то
\[
\begin{aligned}
f_{S_n}(t)
&=
E\exp\left(
it\sum_{j=1}^nX_j
\right)\\
&=
E\prod_{j=1}^ne^{itX_j}\\
&=
\prod_{j=1}^nEe^{itX_j}\\
&=
\prod_{j=1}^nf_{X_j}(t)
\end{aligned}
\]

\end{proofbox}

\begin{proofbox}{свойства 6}

Для любых $t,h\in\mathbb{R}$
\[
\begin{aligned}
|f_X(t+h)-f_X(t)|
&=
\left|
E\left(
e^{i(t+h)X}-e^{itX}
\right)
\right|\\
&\leq
E\left|
e^{itX}
\right|
\left|
e^{ihX}-1
\right|\\
&=
E|e^{ihX}-1|
\end{aligned}
\]

Правая часть не зависит от $t$ и при $h\to0$ стремится к нулю

Следовательно,
\[
\sup_{t\in\mathbb{R}}
|f_X(t+h)-f_X(t)|
\xrightarrow[h\to0]{}0
\]

\end{proofbox}

\begin{example}

\textbf{Характеристические функции некоторых дискретных распределений:}

1) $X\sim I(c) \
\Longrightarrow \
f_X(t)=Ee^{itX}=e^{itc}$

\bigskip

2) Если $X\sim\operatorname{Bern}(p)$, $q=1-p$, то
\[
f_X(t)
=
pe^{it}+qe^{i\cdot0\cdot t}
=
pe^{it}+q
\]

\bigskip

3) Если $X\sim\operatorname{Bin}(n,p)$, то
\[
f_X(t)
=
\left(pe^{it}+q\right)^n
\]

\bigskip

4) Если $X\sim\operatorname{Geom}(p)$, то
\[
\begin{aligned}
f_X(t)
&=
\sum_{k=0}^{\infty}pq^ke^{itk}
=
\frac{p}{1-qe^{it}}
\end{aligned}
\]

\bigskip

5) Если $X\sim\operatorname{NB}(r,p)$, то
\[
f_X(t)
=
\left(
\frac{p}{1-qe^{it}}
\right)^r
\]

\bigskip

6) Если $X\sim\operatorname{Pois}(\lambda)$, то
\[
\begin{aligned}
f_X(t)
&=
\sum_{k=0}^{\infty}
e^{itk}e^{-\lambda}\frac{\lambda^k}{k!}\\
&=
e^{-\lambda}
\sum_{k=0}^{\infty}
\frac{(e^{it}\lambda)^k}{k!}\\
&=
e^{-\lambda+\lambda e^{it}}
=
e^{\lambda(e^{it}-1)}
\end{aligned}
\]

\end{example}

\begin{example}

\textbf{Характеристические функции некоторых абсолютно непрерывных распределений:}

1) Если $X\sim U[a,b]$, то
\[
\begin{aligned}
f_X(t)
&=
\frac{1}{b-a}
\int_a^b e^{itx}\,dx\\
&=
\frac{e^{itb}-e^{ita}}{(b-a)it}
\end{aligned}
\]

При $a=-b$
\[
\begin{aligned}
f_X(t)
&=
\frac{e^{itb}-e^{-itb}}{2bit}\\
&=
\frac{\cos(tb)+i\sin(tb)-\cos(tb)+i\sin(tb)}{2bit}\\
&=
\frac{\sin(tb)}{tb}
\end{aligned}
\]

\bigskip

2) Пусть $X\sim\mathcal{N}(0,1)$. Тогда
\[
f_X(t)
=
Ee^{itX}
=
\frac{1}{\sqrt{2\pi}}
\int_{-\infty}^{+\infty}
e^{itx-\frac{x^2}{2}}\,dx
\]

Дифференцируя по $t$, получаем
\[
\begin{aligned}
f_X'(t)
&=
\frac{1}{\sqrt{2\pi}}
\int_{-\infty}^{+\infty}
ixe^{itx-\frac{x^2}{2}}\,dx\\
&=
-\frac{i}{\sqrt{2\pi}}
\int_{-\infty}^{+\infty}
e^{itx}\,d\left(e^{-\frac{x^2}{2}}\right)\\
&=
-tf_X(t)
\end{aligned}
\]

Причём
\[
f_X(0)=1
\]

Следовательно,
\[
\int\frac{df_X}{dt}
=
-tf
\]
и
\[
\int\frac{df_X}{f_X}
=
\int-t\,dt
\]

Заметим, что
\[
f_X(t)
=
e^{-\frac{t^2}{2}}
\]

\bigskip

3) Для
\[
X\sim\mathcal{N}(\mu,\sigma^2)
\]
имеем
\[
X=\sigma Z+\mu,
\qquad
Z\sim\mathcal{N}(0,1)
\]
поэтому
\[
f_X(t)
=
e^{it\mu}f_Z(\sigma t)
=
e^{it\mu-\frac{\sigma^2t^2}{2}}
\]

\bigskip

4) Если $X\sim\operatorname{Exp}(\lambda)$, то
\[
\begin{aligned}
f_X(t)
&=
\lambda
\int_0^{+\infty}
e^{itx-\lambda x}\,dx\\
&=
\frac{\lambda}{it-\lambda}
\left.e^{(it-\lambda)x}\right|_0^{+\infty}\\
&=
\frac{\lambda}{\lambda-it}
\end{aligned}
\]

\bigskip

5) Если $X\sim\Gamma(\alpha,\lambda)$, то
\[
f_X(t)
=
\int_0^{+\infty}
\frac{\lambda^\alpha x^{\alpha-1}e^{-\lambda x}}
{\Gamma(\alpha)}
e^{itx}\,dx
\]

Далее
\[
\begin{aligned}
f_X'(t)
&=
\int_0^{+\infty}
\frac{\lambda^\alpha x^{\alpha-1}e^{-\lambda x}}
{\Gamma(\alpha)}
ixe^{itx}\,dx
\end{aligned}
\]

Отсюда
\[
f_X(t)
=
\left(
\frac{\lambda}{\lambda-it}
\right)^\alpha
\]

\bigskip

6) Если $X\sim\operatorname{Cauchy}(\mu,\gamma)$, то
\[
f_X(t)
=
e^{i\mu t-\gamma|t|}
\]

\end{example}

\begin{theorem}{о связи распределения и характеристической функции}

Распределение случайной величины однозначно определяет её характеристическую функцию и наоборот

То есть
\[
f_X=f_Y
\quad\Longleftrightarrow\quad
X\overset{d}{=}Y
\]

\end{theorem}

\begin{proofbox}{}

Если
\[
X\overset{d}{=}Y
\]
то непосредственно
\[
f_X(t)
=
Ee^{itX}
=
Ee^{itY}
=
f_Y(t)
\]

Докажем обратное утверждение для целочисленной случайной величины $X$

В этом случае характеристическая функция $2\pi$-периодична:
\[
f_X(0) = f_X(2\pi) = 1
\]

Рассмотрим пространство
\[
L_2[-\pi,\pi]
=
\left\{
f:[-\pi,\pi]\to\mathbb{C}:
\int_{-\pi}^{\pi}|f(t)|^2\,dt<\infty
\right\}
\]

со скалярным произведением
\[
\langle f,g\rangle
=
\int_{-\pi}^{\pi}
f(t)\overline{g(t)}\,dt
\]

Для функций $e^{ikt}$ имеем
\[
\begin{aligned}
\left\langle e^{ikt},e^{ijt}\right\rangle
&=
\int_{-\pi}^{\pi}
e^{i(k-j)t}\,dt\\
&=
\begin{cases}
0, & k\neq j\\
2\pi, & k=j
\end{cases}
\end{aligned}
\]

Поэтому коэффициенты разложения
\[
f_X(t)
=
\sum_k p_ke^{ikt}
\]
однозначно восстанавливаются по формуле
\[
p_k
=
\frac{1}{2\pi}
\int_{-\pi}^{\pi}
f_X(t)e^{-ikt}\,dt
\]

Следовательно, характеристическая функция однозначно определяет вероятности $p_k$, а значит, и распределение $X$

\end{proofbox}

\begin{theorem}{Леви о непрерывности}

Пусть $X_n,X$ --- случайные величины с характеристическими функциями
$f_{X_n}$ и $f_X$

Тогда
\[
X_n\xrightarrow{d}X
\quad\Longleftrightarrow\quad
f_{X_n}(t)\xrightarrow[n\to\infty]{}f_X(t)
\qquad
\forall t\in\mathbb{R}
\]

\end{theorem}

\begin{theorem}{закон больших чисел для независимых одинаково распределённых случайных величин}

Пусть $X_1,X_2,\ldots$ --- независимые одинаково распределённые случайные величины и
\[
EX_k=\mu
\]

Тогда
\[
\frac{1}{n}\sum_{k=1}^nX_k
\xrightarrow{P}
\mu
\]

\end{theorem}

\begin{proofbox}{}

Для характеристической функции $X_k$ при $t\to0$
\[
f_X(t)
=
1+it\mu+o(t)
\]

Обозначим
\[
S_n=\sum_{k=1}^nX_k
\]

Так как случайные величины независимы,
\[
f_{S_n}(t)
=
\left(1+it\mu+o(t)\right)^n
\]

Следовательно,
\[
\begin{aligned}
f_{\frac{S_n}{n}}(t)
=
f_{S_n}\left(\frac{t}{n}\right)
=
\left(
1+\frac{it\mu}{n}
+
o\left(\frac{1}{n}\right)
\right)^n
&\xrightarrow[n\to\infty]{}
e^{it\mu}
\end{aligned}
\]

Но $e^{it\mu}$ --- характеристическая функция константы $\mu$

По теореме Леви
\[
\frac{S_n}{n}
\xrightarrow{d}
\mu
\]

Так как предел является константой,
\[
\frac{S_n}{n}
\xrightarrow{P}
\mu
\]

\end{proofbox}

\begin{theorem}{центральная предельная теорема для независимых одинаково распределённых случайных величин}

Пусть $X_1,X_2,\ldots$ --- независимые одинаково распределённые случайные величины,
\[
EX_k=\mu,
\qquad
\operatorname{Var}X_k=\sigma^2>0
\]

Тогда
\[
\frac{\sum_{k=1}^nX_k-n\mu}
{\sqrt{n\sigma^2}}
\xrightarrow{d}
U\sim\mathcal{N}(0,1)
\]

\end{theorem}

\begin{proofbox}{}

Положим
\[
U_k
=
\frac{X_k-\mu}{\sigma}
\]

Тогда
\[
EU_k=0,
\qquad
\operatorname{Var}U_k=1,
\qquad
EU_k^2=1
\]

При $t\to0$
\[
f_{U_k}(t)
=
1-\frac{t^2}{2}+o(t^2)
\]

Следовательно,
\[
\begin{aligned}
f_{\frac{1}{\sqrt{n}}\sum_{k=1}^nU_k}(t)
=
\left(
f_{U_k}\left(\frac{t}{\sqrt{n}}\right)
\right)^n
=
\left(
1-\frac{t^2}{2n}
+
o\left(\frac{1}{n}\right)
\right)^n
&\xrightarrow[n\to\infty]{}
e^{-\frac{t^2}{2}}
\end{aligned}
\]

Функция
\[
e^{-\frac{t^2}{2}}
\]
является характеристической функцией распределения $\mathcal{N}(0,1)$

По теореме Леви
\[
\frac{\sum_{k=1}^nX_k-n\mu}
{\sqrt{n\sigma^2}}
\xrightarrow{d}
\mathcal{N}(0,1)
\]

\end{proofbox}

\begin{corollary}

При достаточно больших $n$
\[
P\left(
a<
\frac{\sum_{k=1}^nX_k-n\mu}
{\sqrt{n\sigma^2}}
<b
\right)
\approx
\Phi(b)-\Phi(a)
\]

\end{corollary}



\end{document}